% =====================================================================
%  Phương pháp — LCF (CDW / CDM) và Token Merging (3 chiến lược)
%  Dùng % =====================================================================
%  Phương pháp — LCF (CDW / CDM) và Token Merging (3 chiến lược)
%  Dùng % =====================================================================
%  Phương pháp — LCF (CDW / CDM) và Token Merging (3 chiến lược)
%  Dùng % =====================================================================
%  Phương pháp — LCF (CDW / CDM) và Token Merging (3 chiến lược)
%  Dùng \input{phuong_phap_lcf_tome} trong main.tex
% =====================================================================

% =====================================================================
\section{Cơ chế Local Context Focus}
\label{sec:pp_lcf}
% =====================================================================

\subsection{Định dạng đầu vào SPC}
\label{subsec:spc}

Mô hình \texttt{FastLcfBertMultiTask} nhận đầu vào theo định dạng
\textit{Single-Pair Classification} (SPC): cả câu đánh giá lẫn cụm từ
khía cạnh được nối thành một chuỗi duy nhất và đưa vào backbone BERT:

\begin{equation}
  \mathbf{X} = [\texttt{[CLS]},\; x_1, \dots, x_m,\;
                \texttt{[SEP]},\; a_1, \dots, a_k,\;
                \texttt{[SEP]}]
  \label{eq:spc}
\end{equation}

\noindent trong đó $x_1, \dots, x_m$ là các token của câu văn và
$a_1, \dots, a_k$ là các token của cụm từ khía cạnh.  Tổng độ dài
chuỗi là $L = m + k + 3$ (bao gồm hai token \texttt{[SEP]} và một
\texttt{[CLS]}).

Bên cạnh đó, một \textbf{véc-tơ chỉ thị khía cạnh} (aspect indicator)
$\mathbf{v} \in \{0, 1\}^{L}$ được xây dựng tương ứng:

\begin{equation}
  v_i = \begin{cases}
    1 & \text{nếu token thứ } i \text{ thuộc vị trí cụm khía cạnh
        (segment A, khớp với span gốc)}\\
    0 & \text{ngược lại}
  \end{cases}
  \label{eq:aspect_indicator}
\end{equation}

Véc-tơ $\mathbf{v}$ đóng vai trò đầu vào cho tất cả các cơ chế LCF và
được truyền xuyên suốt qua các module Token Merging (mục~\ref{sec:pp_tome}).

% ---------------------------------------------------------------------
\subsection{Khoảng cách ngữ nghĩa tương đối (SRD)}
\label{subsec:srd}
% ---------------------------------------------------------------------

Để định lượng mức độ ``gần'' của một token $i$ so với cụm từ khía
cạnh, luận văn sử dụng \textbf{Semantic Relative Distance} (SRD) theo
định nghĩa trong LCF-ATEPC~\cite{zeng2019lcf}:

\begin{equation}
  \mathrm{SRD}_i = \max\!\left(0,\;
    \left|i - c\right| - \left\lfloor \frac{l_a}{2} \right\rfloor
  \right)
  \label{eq:srd}
\end{equation}

\noindent với:

\begin{itemize}
  \item $c = \dfrac{a_{\mathrm{start}} + a_{\mathrm{end}}}{2}$: chỉ số
    trung tâm của cụm từ khía cạnh trong chuỗi token;
  \item $l_a = a_{\mathrm{end}} - a_{\mathrm{start}} + 1$: độ dài (số
    token) của cụm từ khía cạnh;
  \item $a_{\mathrm{start}}, a_{\mathrm{end}}$: chỉ số token đầu và
    cuối của cụm khía cạnh trong chuỗi đầy đủ.
\end{itemize}

Khi token $i$ nằm trong vùng trực tiếp bao phủ bởi cụm từ khía cạnh,
$\mathrm{SRD}_i = 0$.  Khi token nằm ngoài vùng đó, $\mathrm{SRD}_i$
tăng tuyến tính theo khoảng cách từ biên cụm khía cạnh đến token $i$.

Ngưỡng bán kính $\alpha$ (tham số \texttt{srd\_threshold}, mặc định
$\alpha = 5$) xác định vùng ngữ cảnh cục bộ: mọi token có
$\mathrm{SRD}_i \leq \alpha$ được coi là \textit{nằm trong ngữ cảnh
cục bộ} và nhận trọng số tối đa bằng 1.

% ---------------------------------------------------------------------
\subsection{Context Dynamic Weighting (CDW)}
\label{subsec:cdw}
% ---------------------------------------------------------------------

CDW gán cho mỗi token một \textbf{trọng số liên tục} $w_i \in [0, 1]$
phản ánh mức độ liên quan đến khía cạnh đang xét:

\begin{equation}
  w_i^{\mathrm{CDW}} = \begin{cases}
    1 & \text{nếu } \mathrm{SRD}_i \leq \alpha \\[6pt]
    \dfrac{L - (\mathrm{SRD}_i - \alpha)}{L}
      & \text{nếu } \mathrm{SRD}_i > \alpha
  \end{cases}
  \label{eq:cdw}
\end{equation}

\noindent sau đó được kẹp vào đoạn $[0, 1]$:

\begin{equation}
  w_i^{\mathrm{CDW}} \leftarrow \mathrm{clamp}\!\left(
    w_i^{\mathrm{CDW}},\; 0,\; 1
  \right)
  \label{eq:cdw_clamp}
\end{equation}

Vùng trong bán kính $\alpha$ nhận trọng số 1 (không suy giảm); càng
ra xa khía cạnh, trọng số giảm tuyến tính và bằng 0 khi
$\mathrm{SRD}_i = \alpha + L$.  Trong trường hợp không có token khía
cạnh nào được phát hiện trong câu (véc-tơ $\mathbf{v}$ toàn 0), CDW
trả về véc-tơ trọng số hằng số bằng 1 cho toàn bộ chuỗi.

Véc-tơ trọng số CDW $\mathbf{w}^{\mathrm{CDW}} \in \mathbb{R}^L$ được
nhân phần tử với ma trận trạng thái ẩn của BERT để tạo ra
\textbf{luồng cục bộ} (local stream):

\begin{equation}
  \mathbf{H}^{\mathrm{local}} = \mathbf{H} \odot
    \mathbf{w}^{\mathrm{CDW}} \mathbf{1}_{H}^{\top}
  \label{eq:local_cdw}
\end{equation}

\noindent trong đó $\mathbf{H} \in \mathbb{R}^{L \times H}$ là đầu ra
của backbone BERT, $\odot$ là tích Hadamard theo chiều chuỗi (mỗi
hàng $\mathbf{H}_i$ được nhân với vô hướng $w_i$), và
$\mathbf{1}_{H}^{\top}$ ký hiệu phép broadcast sang chiều ẩn $H$.

% ---------------------------------------------------------------------
\subsection{Context Dynamic Masking (CDM)}
\label{subsec:cdm}
% ---------------------------------------------------------------------

CDM là biến thể nhị phân của CDW: thay vì giảm trọng số liên tục, nó
\textbf{triệt tiêu hoàn toàn} những token nằm ngoài ngưỡng $\alpha$:

\begin{equation}
  w_i^{\mathrm{CDM}} = \begin{cases}
    1 & \text{nếu } \mathrm{SRD}_i \leq \alpha \\[4pt]
    0 & \text{nếu } \mathrm{SRD}_i > \alpha
  \end{cases}
  \label{eq:cdm}
\end{equation}

So với CDW, mặt nạ CDM tạo ra một ranh giới cứng quanh vùng ngữ cảnh
cục bộ. Luồng cục bộ khi dùng CDM là:

\begin{equation}
  \mathbf{H}^{\mathrm{local}} = \mathbf{H} \odot
    \mathbf{w}^{\mathrm{CDM}} \mathbf{1}_{H}^{\top}
  \label{eq:local_cdm}
\end{equation}

Khi không có token khía cạnh, CDM cũng trả về véc-tơ trọng số toàn 1
(tương tự CDW), đảm bảo mô hình không suy sụp hoàn toàn trên các mẫu
thiếu annotation.

% ---------------------------------------------------------------------
\subsection{Kiến trúc hai luồng và phân loại đa nhiệm}
\label{subsec:two_stream}
% ---------------------------------------------------------------------

Sau khi tính luồng cục bộ bằng CDW hoặc CDM, mô hình thực hiện
\textbf{thiết kế hai luồng} (two-stream) kết hợp thông tin cục bộ
và toàn cục:

\paragraph{Bước 1: Self-attention cục bộ.}
Luồng cục bộ được đưa qua một lớp self-attention bổ sung
$\mathrm{SA}_{\mathrm{LCF}}$ với kết nối dư (residual) và
chuẩn hóa lớp (LayerNorm):

\begin{equation}
  \widetilde{\mathbf{H}}^{\mathrm{local}} =
    \mathrm{LayerNorm}\!\left(
      \mathbf{H}^{\mathrm{local}} +
      \mathrm{Dropout}\!\left(
        \mathrm{MHA}\!\left(
          \mathbf{H}^{\mathrm{local}},\,
          \mathbf{H}^{\mathrm{local}},\,
          \mathbf{H}^{\mathrm{local}}
        \right)
      \right)
    \right)
  \label{eq:sa_lcf}
\end{equation}

\noindent trong đó $\mathrm{MHA}(\cdot)$ là phép multi-head attention
với 8 head.

\paragraph{Bước 2: Hợp nhất hai luồng.}
Luồng cục bộ và luồng toàn cục được ghép nối theo chiều ẩn rồi chiếu
xuống $H$ chiều:

\begin{equation}
  \mathbf{H}^{\mathrm{fused}} = \mathrm{Dropout}\!\left(
    \mathbf{W}_2 \left[
      \widetilde{\mathbf{H}}^{\mathrm{local}};
      \mathbf{H}
    \right]
  \right), \quad \mathbf{W}_2 \in \mathbb{R}^{H \times 2H}
  \label{eq:fusion}
\end{equation}

\paragraph{Bước 3: Self-attention tổng hợp và pooling.}
Biểu diễn tổng hợp tiếp tục qua một lớp self-attention $\mathrm{SA}$
nữa trước khi lấy biểu diễn token \texttt{[CLS]}:

\begin{equation}
  \mathbf{H}^{\mathrm{out}} = \mathrm{SA}\!\left(
    \mathbf{H}^{\mathrm{fused}}
  \right)
  \label{eq:sa_final}
\end{equation}

\begin{equation}
  \mathbf{p} = \mathrm{Tanh}\!\left(
    \mathbf{W}_p\, \mathbf{H}^{\mathrm{out}}_{[0]}
  \right), \quad \mathbf{W}_p \in \mathbb{R}^{H \times H}
  \label{eq:pooler}
\end{equation}

\noindent trong đó $\mathbf{H}^{\mathrm{out}}_{[0]}$ là hàng ứng với
vị trí \texttt{[CLS]} (chỉ số 0).

\paragraph{Bước 4: Phân loại đa nhiệm.}
Từ biểu diễn câu $\mathbf{p} \in \mathbb{R}^H$, hai đầu ra được tính
song song:

\begin{align}
  \hat{\mathbf{y}}^{\mathrm{sent}} &= \mathbf{W}_s\, \mathbf{p},
    \quad \mathbf{W}_s \in \mathbb{R}^{3 \times H}
    \label{eq:head_sentiment}\\[4pt]
  \hat{\mathbf{y}}^{\mathrm{cat}} &= \mathbf{W}_c\, \mathbf{p},
    \quad \mathbf{W}_c \in \mathbb{R}^{6 \times H}
    \label{eq:head_category}
\end{align}

\noindent với $\hat{\mathbf{y}}^{\mathrm{sent}}$ là logit cho 3 lớp
cảm xúc (Positive / Negative / Neutral) và
$\hat{\mathbf{y}}^{\mathrm{cat}}$ là logit cho 6 nhóm khía cạnh
(SERVICE, FACILITY, AMENITY, EXPERIENCE, BRANDING, LOYALTY).

Hàm mất mát tổng hợp là tổng của hai hàm cross-entropy có trọng số
lớp:

\begin{equation}
  \mathcal{L} = \mathcal{L}_{\mathrm{CE}}^{\mathrm{sent}}
              + \mathcal{L}_{\mathrm{CE}}^{\mathrm{cat}}
  \label{eq:loss}
\end{equation}

Hình~\ref{fig:two_stream} minh họa toàn bộ kiến trúc hai luồng.

% =====================================================================
\section{Token Merging cho chuỗi 1D}
\label{sec:pp_tome}
% =====================================================================

\subsection{Tổng quan và vị trí tích hợp}
\label{subsec:tome_overview}

Token Merging (ToMe)~\cite{bolya2023tome} được tích hợp vào mô hình
\texttt{FastLcfBertMultiTask} dưới dạng module
\texttt{ToMeSequenceMerger}, áp dụng \textbf{sau} khi backbone BERT
hoàn thành quá trình mã hóa (post-BERT ToMe) và trước khi luồng cục bộ
được tính.

Module nhận đầu vào ba tensor:
$\mathbf{H} \in \mathbb{R}^{B \times L \times H}$ (trạng thái ẩn),
$\mathbf{V} \in \mathbb{R}^{B \times L}$ (véc-tơ chỉ thị khía cạnh)
và $\mathbf{M} \in \{0,1\}^{B \times L}$ (attention mask),
và trả về bộ $(\mathbf{H}', \mathbf{V}', \mathbf{M}')$ có độ dài
chuỗi $L' \leq L$.

Số vòng gộp (merge steps) được cấu hình qua
\texttt{num\_merge\_steps}$= 2$ (mặc định); mỗi vòng áp dụng một
trong ba chiến lược sau.

% ---------------------------------------------------------------------
\subsection{Phép gộp cơ bản và bảo toàn LCF}
\label{subsec:merge_op}
% ---------------------------------------------------------------------

Bất kể chiến lược nào, khi hai token $s$ (source) và $d$ (destination)
được chọn để gộp, phép gộp thực hiện hai thao tác:

\paragraph{Cập nhật biểu diễn bằng trung bình:}
\begin{equation}
  \mathbf{h}'_s = \frac{\mathbf{h}_s + \mathbf{h}_d}{2}
  \label{eq:merge_avg}
\end{equation}

\paragraph{Bảo toàn thông tin khía cạnh bằng max:}
\begin{equation}
  v'_s = \max(v_s,\; v_d)
  \label{eq:merge_lcf}
\end{equation}

Token $d$ bị loại khỏi chuỗi; token $s$ thay thế với giá trị đã được
cập nhật.  Quy tắc \eqref{eq:merge_lcf} đảm bảo rằng nếu một trong
hai token thuộc cụm từ khía cạnh (có $v = 1$), token hợp nhất kết quả
cũng mang nhãn khía cạnh, tránh mất tín hiệu vị trí cho cơ chế LCF
ở bước tiếp theo.

Tất cả ba chiến lược dùng \textbf{độ tương đồng cosine} trên biểu
diễn đã chuẩn hóa $L_2$ để đo mức độ giống nhau giữa hai token:

\begin{equation}
  \hat{\mathbf{h}}_i = \frac{\mathbf{h}_i}{\|\mathbf{h}_i\|_2 + \varepsilon},
  \quad \varepsilon = 10^{-6}
  \label{eq:l2norm}
\end{equation}

\begin{equation}
  \mathrm{sim}(i, j) = \hat{\mathbf{h}}_i \cdot \hat{\mathbf{h}}_j
  \label{eq:cosine}
\end{equation}

% ---------------------------------------------------------------------
\subsection{Chiến lược 1: Bipartite Matching}
\label{subsec:bipartite}
% ---------------------------------------------------------------------

Chiến lược này là chiến lược gốc của ToMe~\cite{bolya2023tome}, được
điều chỉnh cho chuỗi 1D và bổ sung bảo vệ token khía cạnh.

\paragraph{Xây dựng tập ứng viên.}
Từ $n$ token hợp lệ của câu (không tính padding), tập ứng viên gộp
$\mathcal{P}$ được xác định bằng cách loại trừ:
\begin{itemize}
  \item Token \texttt{[CLS]} ($i = 0$), bảo vệ bởi \texttt{protect\_cls};
  \item Token \texttt{[SEP]} cuối ($i = n-1$), bảo vệ bởi
    \texttt{protect\_sep};
  \item Các token khía cạnh ($v_i > 0.5$), bảo vệ bởi
    \texttt{protect\_aspect}.
\end{itemize}

\paragraph{Phân chia hai nhóm.}
Các phần tử của $\mathcal{P}$ được liệt kê theo thứ tự chỉ số tăng
dần thành dãy $p_0, p_1, \dots, p_{|\mathcal{P}|-1}$.  Dãy này được
chia xen kẽ thành hai nhóm:

\begin{align}
  \mathcal{A} &= \{p_0,\, p_2,\, p_4, \dots\} \quad \text{(chỉ số
    chẵn trong dãy)}\label{eq:group_A}\\
  \mathcal{B} &= \{p_1,\, p_3,\, p_5, \dots\} \quad \text{(chỉ số
    lẻ trong dãy)}\label{eq:group_B}
\end{align}

\paragraph{Tính ma trận độ tương đồng.}
Mỗi hàng ứng với một token trong $\mathcal{A}$, mỗi cột ứng với một
token trong $\mathcal{B}$:

\begin{equation}
  \mathbf{S} = \widehat{\mathbf{A}}\, \widehat{\mathbf{B}}^{\top},
  \quad \mathbf{S} \in \mathbb{R}^{|\mathcal{A}| \times |\mathcal{B}|}
  \label{eq:bipartite_score}
\end{equation}

\noindent trong đó $\widehat{\mathbf{A}}, \widehat{\mathbf{B}}$ là ma
trận biểu diễn chuẩn hóa $L_2$ của các token trong nhóm tương ứng.

\paragraph{Ghép cặp tham lam.}
Mỗi token $a_i \in \mathcal{A}$ chọn token $b_{j^*} \in \mathcal{B}$
có độ tương đồng lớn nhất:

\begin{equation}
  j^*(i) = \arg\max_{j} S_{ij}
  \label{eq:greedy_assign}
\end{equation}

Ràng buộc một-một phía $\mathcal{B}$: mỗi $b_j$ chỉ được gộp với một
$a_i$ duy nhất (cặp đầu tiên yêu cầu $b_j$ được giữ lại, các cặp sau
yêu cầu cùng $b_j$ bị bỏ qua).  Đây là phép ghép cặp hai phía
\textit{greedy} theo nghĩa mỗi bên $\mathcal{A}$ giải quyết độc lập
mà không có giai đoạn đấu giá toàn cục.

Sau bước này, mỗi cặp $(a_i, b_{j^*(i)})$ được gộp theo
\eqref{eq:merge_avg}--\eqref{eq:merge_lcf} với $s = a_i$, $d =
b_{j^*(i)}$.

\paragraph{Độ phức tạp.}
Với $|\mathcal{P}|$ token ứng viên và chiều ẩn $H$, chi phí tính ma
trận $\mathbf{S}$ là $\mathcal{O}(|\mathcal{P}|^2 / 4 \cdot H)$, thấp
hơn nhiều so với self-attention $\mathcal{O}(L^2 H)$ trên toàn chuỗi.

% ---------------------------------------------------------------------
\subsection{Chiến lược 2: Sequential Local Merging}
\label{subsec:sequential}
% ---------------------------------------------------------------------

Chiến lược này thực hiện gộp \textbf{tuần tự từ trái sang phải}, mỗi
token quyết định gộp dựa trên so sánh với hai lân cận gần nhất.

\paragraph{Xây dựng tập bảo vệ.}
Ngoài CLS, SEP và token khía cạnh, chiến lược này bổ sung bảo vệ
\textbf{ranh giới mid-SEP}: trong định dạng SPC
(xem~\eqref{eq:spc}), token phân cách giữa phần câu văn và phần
khía cạnh (tức \texttt{[SEP]} đầu tiên, không phải token cuối) không
được xóa.  Điều kiện nhận biết:

\begin{equation}
  \mathrm{midSep}(j) = \mathbf{1}\!\left[
    v_j < 0.5 \;\text{ và }\; v_{j+1} > 0.5
  \right]
  \label{eq:midsep}
\end{equation}

\paragraph{Vòng quét.}
Biến con trỏ $i$ duyệt từ $\texttt{protect\_left}$ đến $n -
\texttt{protect\_right} - 1$.  Tại mỗi bước $i$:

\begin{enumerate}
  \item Bỏ qua nếu $i$ đã bị xóa, là mid-SEP hoặc là token khía
    cạnh.
  \item Tìm \textbf{lân cận trái tích cực gần nhất}:
    $l = \max\{l' < i : l' \text{ chưa xóa và không là mid-SEP}\}$.
  \item Tìm \textbf{lân cận phải tích cực gần nhất không phải
    khía cạnh}:
    $r = \min\{r' > i : r' \text{ chưa xóa, không mid-SEP, không
    aspect}\}$.
  \item Tính độ tương đồng:
    \begin{align}
      \mathrm{sim}_{\mathrm{L}} &=
        \mathrm{sim}(i, l) \quad \text{(nếu tồn tại } l \text{)}
        \label{eq:sim_left}\\
      \mathrm{sim}_{\mathrm{R}} &=
        \mathrm{sim}(i, r) \quad \text{(nếu tồn tại } r \text{)}
        \label{eq:sim_right}
    \end{align}
  \item Ra quyết định gộp:
    \begin{equation}
      \begin{cases}
        \text{Token } i \text{ gộp vào trái}:
          & \mathbf{h}_l \leftarrow \tfrac{\mathbf{h}_l + \mathbf{h}_i}{2},\;
            v_l \leftarrow \max(v_l, v_i),\;
            \text{xóa } i \\
            & \quad \text{khi } \exists l \text{ và }
              \mathrm{sim}_{\mathrm{L}} > \mathrm{sim}_{\mathrm{R}}\\[6pt]
        \text{Token phải gộp vào } i:
          & \mathbf{h}_i \leftarrow \tfrac{\mathbf{h}_i + \mathbf{h}_r}{2},\;
            v_i \leftarrow \max(v_i, v_r),\;
            \text{xóa } r \\
            & \quad \text{khi } \exists r \text{ và }
              \mathrm{sim}_{\mathrm{R}} \geq \mathrm{sim}_{\mathrm{L}}
      \end{cases}
      \label{eq:seq_decision}
    \end{equation}
  \item Tăng $i \leftarrow i + 1$.
\end{enumerate}

Một lần duyệt đầy đủ ($i$ từ đầu đến cuối) được gọi là một
\textit{merge step}; mô hình thực hiện \texttt{num\_merge\_steps}$= 2$
lần duyệt liên tiếp.

\paragraph{Tính chất.}
Token có thể \textit{nhận} nhiều gộp từ nhiều hướng trong cùng một
lần duyệt (không có ràng buộc một-một phía đích).  Điều này khác với
Bipartite Matching và cho phép co chuỗi nhanh hơn khi có nhiều cụm
token lặp lại liên tiếp.

% ---------------------------------------------------------------------
\subsection{Chiến lược 3: Attention-Weighted Merging}
\label{subsec:attn_merge}
% ---------------------------------------------------------------------

Chiến lược này gộp ưu tiên \textbf{các token có trọng số chú ý thấp
nhất}, đồng thời bảo vệ các token mang thông tin quan trọng.

\paragraph{Xây dựng tập bảo vệ cứng.}
Ba nhóm token không bao giờ bị xóa:

\begin{itemize}
  \item \textbf{Biên cấu trúc}: CLS ($i = 0$) và SEP ($i = n-1$);
  \item \textbf{Token khía cạnh}: $v_i > 0.5$;
  \item \textbf{Token chú ý cao}: những token có trọng số chú ý
    $a_i$ nằm trong \textbf{phân vị thứ 75} trở lên:
    \begin{equation}
      \theta = \mathrm{quantile}(\mathbf{a},\; 0.75)
      \label{eq:attn_threshold}
    \end{equation}
    \begin{equation}
      \mathrm{highAttn}(i) = \mathbf{1}[a_i \geq \theta]
      \label{eq:high_attn}
    \end{equation}
\end{itemize}

Token $i$ bị bảo vệ khi thỏa ít nhất một trong các điều kiện trên.

\paragraph{Giới hạn số cặp gộp.}
\begin{equation}
  T_{\max} = \left\lfloor
    \frac{n - \texttt{protect\_left} - \texttt{protect\_right}}{2}
  \right\rfloor
  \label{eq:max_merges}
\end{equation}

\paragraph{Vòng lặp gộp lặp.}
Trong mỗi bước $t = 1, 2, \dots, T_{\max}$:

\begin{enumerate}
  \item \textbf{Tìm ứng viên bị xóa}: token chỉ số $i$ có trọng số
    chú ý nhỏ nhất trong số các token chưa bị xóa và không thuộc tập
    bảo vệ:
    \begin{equation}
      i = \arg\min_{\substack{t:\, \text{not removed}(t)\\
            \text{not protected}(t),\, v_t \leq 0.5}}
        a_t
      \label{eq:find_min_attn}
    \end{equation}
    Nếu không có ứng viên: dừng vòng lặp sớm.

  \item \textbf{Tìm token đích}: token $j \neq i$ chưa bị xóa, không
    phải khía cạnh, có độ tương đồng cosine với $i$ cao nhất:
    \begin{equation}
      j = \arg\max_{\substack{k \neq i:\, \text{not removed}(k)\\
            v_k \leq 0.5}}
        \hat{\mathbf{h}}_i \cdot \hat{\mathbf{h}}_k
      \label{eq:find_neighbor}
    \end{equation}
    Nếu giá trị tương đồng $< -1.5$ (không có lân cận hợp lệ): dừng.

  \item \textbf{Thực hiện gộp}: gộp $i$ vào $j$ theo
    \eqref{eq:merge_avg}--\eqref{eq:merge_lcf}, đánh dấu $i$ đã xóa.

  \item \textbf{Cập nhật metric}: chuẩn hóa lại toàn bộ ma trận biểu
    diễn trước vòng lặp kế tiếp.
\end{enumerate}

\paragraph{Sự khác biệt chính so với hai chiến lược còn lại.}
Bipartite Matching và Sequential Local chọn cặp dựa hoàn toàn vào độ
tương đồng cosine giữa biểu diễn.  Attention-Weighted Merging sử dụng
trọng số chú ý $\mathbf{a}$ như một \textit{tín hiệu tầm quan trọng}
độc lập để ưu tiên xóa các token mang ít thông tin cảm xúc nhất, đồng
thời bảo vệ không chỉ token khía cạnh mà cả những token có cơ chế chú
ý của BERT đã học được rằng chúng quan trọng.

% ---------------------------------------------------------------------
\subsection{Chế độ resize và compact}
\label{subsec:resize}
% ---------------------------------------------------------------------

Sau tất cả các vòng gộp, chuỗi có độ dài $L' \leq L$.  Module
\texttt{ToMeSequenceMerger} hỗ trợ hai chế độ đầu ra, điều khiển bởi
tham số \texttt{resize}:

\paragraph{Chế độ \texttt{resize=True} (nội suy về độ dài gốc).}
Chuỗi $\mathbf{H}' \in \mathbb{R}^{L' \times H}$ được nội suy tuyến
tính về đúng độ dài $L$ ban đầu:

\begin{equation}
  \mathbf{H}^{\mathrm{out}} =
    \mathrm{Interpolate}\!\left(\mathbf{H}',\; L,\;
    \text{mode=linear}\right)
    \in \mathbb{R}^{L \times H}
  \label{eq:resize}
\end{equation}

Chế độ này \textbf{không thay đổi hình dạng tensor} nên các lớp hạ
nguồn (SA, Linear) không cần điều chỉnh.  Nó cô lập tác động của
gộp token lên \textit{chất lượng biểu diễn} mà không đo tốc độ suy
luận.  Đây là chế độ được dùng trong toàn bộ các thực nghiệm chính của
luận văn.

\paragraph{Chế độ \texttt{resize=False} (compact).}
Chuỗi rút gọn $\mathbf{H}'$ được giữ nguyên; attention mask mới
$\mathbf{M}' \in \{0, 1\}^{B \times L'_{\max}}$ được tính và truyền
xuống các lớp self-attention.  Chế độ này mang lại \textbf{tăng tốc
thực sự} ($\mathcal{O}({L'}^2)$ thay cho $\mathcal{O}(L^2)$) nhưng
yêu cầu các lớp hạ nguồn hỗ trợ chuỗi có độ dài thay đổi.  Đây là
hướng tối ưu hóa được đề xuất cho các nghiên cứu tiếp theo.

\paragraph{Tóm tắt hai chế độ.}

\begin{table}[h]
  \centering
  \caption{So sánh hai chế độ đầu ra của \texttt{ToMeSequenceMerger}.}
  \label{tab:resize_mode}
  \begin{tabular}{lcc}
    \toprule
    \textbf{Thuộc tính} & \texttt{resize=True} & \texttt{resize=False} \\
    \midrule
    Độ dài đầu ra & $L$ (giữ nguyên) & $L' \leq L$ \\
    Hình dạng tensor thay đổi & Không & Có \\
    Tăng tốc suy luận thực sự & Không & Có \\
    Mục tiêu đánh giá & Chất lượng biểu diễn & Tốc độ + chất lượng \\
    Dùng trong thực nghiệm chính & \checkmark & \\
    \bottomrule
  \end{tabular}
\end{table}
 trong main.tex
% =====================================================================

% =====================================================================
\section{Cơ chế Local Context Focus}
\label{sec:pp_lcf}
% =====================================================================

\subsection{Định dạng đầu vào SPC}
\label{subsec:spc}

Mô hình \texttt{FastLcfBertMultiTask} nhận đầu vào theo định dạng
\textit{Single-Pair Classification} (SPC): cả câu đánh giá lẫn cụm từ
khía cạnh được nối thành một chuỗi duy nhất và đưa vào backbone BERT:

\begin{equation}
  \mathbf{X} = [\texttt{[CLS]},\; x_1, \dots, x_m,\;
                \texttt{[SEP]},\; a_1, \dots, a_k,\;
                \texttt{[SEP]}]
  \label{eq:spc}
\end{equation}

\noindent trong đó $x_1, \dots, x_m$ là các token của câu văn và
$a_1, \dots, a_k$ là các token của cụm từ khía cạnh.  Tổng độ dài
chuỗi là $L = m + k + 3$ (bao gồm hai token \texttt{[SEP]} và một
\texttt{[CLS]}).

Bên cạnh đó, một \textbf{véc-tơ chỉ thị khía cạnh} (aspect indicator)
$\mathbf{v} \in \{0, 1\}^{L}$ được xây dựng tương ứng:

\begin{equation}
  v_i = \begin{cases}
    1 & \text{nếu token thứ } i \text{ thuộc vị trí cụm khía cạnh
        (segment A, khớp với span gốc)}\\
    0 & \text{ngược lại}
  \end{cases}
  \label{eq:aspect_indicator}
\end{equation}

Véc-tơ $\mathbf{v}$ đóng vai trò đầu vào cho tất cả các cơ chế LCF và
được truyền xuyên suốt qua các module Token Merging (mục~\ref{sec:pp_tome}).

% ---------------------------------------------------------------------
\subsection{Khoảng cách ngữ nghĩa tương đối (SRD)}
\label{subsec:srd}
% ---------------------------------------------------------------------

Để định lượng mức độ ``gần'' của một token $i$ so với cụm từ khía
cạnh, luận văn sử dụng \textbf{Semantic Relative Distance} (SRD) theo
định nghĩa trong LCF-ATEPC~\cite{zeng2019lcf}:

\begin{equation}
  \mathrm{SRD}_i = \max\!\left(0,\;
    \left|i - c\right| - \left\lfloor \frac{l_a}{2} \right\rfloor
  \right)
  \label{eq:srd}
\end{equation}

\noindent với:

\begin{itemize}
  \item $c = \dfrac{a_{\mathrm{start}} + a_{\mathrm{end}}}{2}$: chỉ số
    trung tâm của cụm từ khía cạnh trong chuỗi token;
  \item $l_a = a_{\mathrm{end}} - a_{\mathrm{start}} + 1$: độ dài (số
    token) của cụm từ khía cạnh;
  \item $a_{\mathrm{start}}, a_{\mathrm{end}}$: chỉ số token đầu và
    cuối của cụm khía cạnh trong chuỗi đầy đủ.
\end{itemize}

Khi token $i$ nằm trong vùng trực tiếp bao phủ bởi cụm từ khía cạnh,
$\mathrm{SRD}_i = 0$.  Khi token nằm ngoài vùng đó, $\mathrm{SRD}_i$
tăng tuyến tính theo khoảng cách từ biên cụm khía cạnh đến token $i$.

Ngưỡng bán kính $\alpha$ (tham số \texttt{srd\_threshold}, mặc định
$\alpha = 5$) xác định vùng ngữ cảnh cục bộ: mọi token có
$\mathrm{SRD}_i \leq \alpha$ được coi là \textit{nằm trong ngữ cảnh
cục bộ} và nhận trọng số tối đa bằng 1.

% ---------------------------------------------------------------------
\subsection{Context Dynamic Weighting (CDW)}
\label{subsec:cdw}
% ---------------------------------------------------------------------

CDW gán cho mỗi token một \textbf{trọng số liên tục} $w_i \in [0, 1]$
phản ánh mức độ liên quan đến khía cạnh đang xét:

\begin{equation}
  w_i^{\mathrm{CDW}} = \begin{cases}
    1 & \text{nếu } \mathrm{SRD}_i \leq \alpha \\[6pt]
    \dfrac{L - (\mathrm{SRD}_i - \alpha)}{L}
      & \text{nếu } \mathrm{SRD}_i > \alpha
  \end{cases}
  \label{eq:cdw}
\end{equation}

\noindent sau đó được kẹp vào đoạn $[0, 1]$:

\begin{equation}
  w_i^{\mathrm{CDW}} \leftarrow \mathrm{clamp}\!\left(
    w_i^{\mathrm{CDW}},\; 0,\; 1
  \right)
  \label{eq:cdw_clamp}
\end{equation}

Vùng trong bán kính $\alpha$ nhận trọng số 1 (không suy giảm); càng
ra xa khía cạnh, trọng số giảm tuyến tính và bằng 0 khi
$\mathrm{SRD}_i = \alpha + L$.  Trong trường hợp không có token khía
cạnh nào được phát hiện trong câu (véc-tơ $\mathbf{v}$ toàn 0), CDW
trả về véc-tơ trọng số hằng số bằng 1 cho toàn bộ chuỗi.

Véc-tơ trọng số CDW $\mathbf{w}^{\mathrm{CDW}} \in \mathbb{R}^L$ được
nhân phần tử với ma trận trạng thái ẩn của BERT để tạo ra
\textbf{luồng cục bộ} (local stream):

\begin{equation}
  \mathbf{H}^{\mathrm{local}} = \mathbf{H} \odot
    \mathbf{w}^{\mathrm{CDW}} \mathbf{1}_{H}^{\top}
  \label{eq:local_cdw}
\end{equation}

\noindent trong đó $\mathbf{H} \in \mathbb{R}^{L \times H}$ là đầu ra
của backbone BERT, $\odot$ là tích Hadamard theo chiều chuỗi (mỗi
hàng $\mathbf{H}_i$ được nhân với vô hướng $w_i$), và
$\mathbf{1}_{H}^{\top}$ ký hiệu phép broadcast sang chiều ẩn $H$.

% ---------------------------------------------------------------------
\subsection{Context Dynamic Masking (CDM)}
\label{subsec:cdm}
% ---------------------------------------------------------------------

CDM là biến thể nhị phân của CDW: thay vì giảm trọng số liên tục, nó
\textbf{triệt tiêu hoàn toàn} những token nằm ngoài ngưỡng $\alpha$:

\begin{equation}
  w_i^{\mathrm{CDM}} = \begin{cases}
    1 & \text{nếu } \mathrm{SRD}_i \leq \alpha \\[4pt]
    0 & \text{nếu } \mathrm{SRD}_i > \alpha
  \end{cases}
  \label{eq:cdm}
\end{equation}

So với CDW, mặt nạ CDM tạo ra một ranh giới cứng quanh vùng ngữ cảnh
cục bộ. Luồng cục bộ khi dùng CDM là:

\begin{equation}
  \mathbf{H}^{\mathrm{local}} = \mathbf{H} \odot
    \mathbf{w}^{\mathrm{CDM}} \mathbf{1}_{H}^{\top}
  \label{eq:local_cdm}
\end{equation}

Khi không có token khía cạnh, CDM cũng trả về véc-tơ trọng số toàn 1
(tương tự CDW), đảm bảo mô hình không suy sụp hoàn toàn trên các mẫu
thiếu annotation.

% ---------------------------------------------------------------------
\subsection{Kiến trúc hai luồng và phân loại đa nhiệm}
\label{subsec:two_stream}
% ---------------------------------------------------------------------

Sau khi tính luồng cục bộ bằng CDW hoặc CDM, mô hình thực hiện
\textbf{thiết kế hai luồng} (two-stream) kết hợp thông tin cục bộ
và toàn cục:

\paragraph{Bước 1: Self-attention cục bộ.}
Luồng cục bộ được đưa qua một lớp self-attention bổ sung
$\mathrm{SA}_{\mathrm{LCF}}$ với kết nối dư (residual) và
chuẩn hóa lớp (LayerNorm):

\begin{equation}
  \widetilde{\mathbf{H}}^{\mathrm{local}} =
    \mathrm{LayerNorm}\!\left(
      \mathbf{H}^{\mathrm{local}} +
      \mathrm{Dropout}\!\left(
        \mathrm{MHA}\!\left(
          \mathbf{H}^{\mathrm{local}},\,
          \mathbf{H}^{\mathrm{local}},\,
          \mathbf{H}^{\mathrm{local}}
        \right)
      \right)
    \right)
  \label{eq:sa_lcf}
\end{equation}

\noindent trong đó $\mathrm{MHA}(\cdot)$ là phép multi-head attention
với 8 head.

\paragraph{Bước 2: Hợp nhất hai luồng.}
Luồng cục bộ và luồng toàn cục được ghép nối theo chiều ẩn rồi chiếu
xuống $H$ chiều:

\begin{equation}
  \mathbf{H}^{\mathrm{fused}} = \mathrm{Dropout}\!\left(
    \mathbf{W}_2 \left[
      \widetilde{\mathbf{H}}^{\mathrm{local}};
      \mathbf{H}
    \right]
  \right), \quad \mathbf{W}_2 \in \mathbb{R}^{H \times 2H}
  \label{eq:fusion}
\end{equation}

\paragraph{Bước 3: Self-attention tổng hợp và pooling.}
Biểu diễn tổng hợp tiếp tục qua một lớp self-attention $\mathrm{SA}$
nữa trước khi lấy biểu diễn token \texttt{[CLS]}:

\begin{equation}
  \mathbf{H}^{\mathrm{out}} = \mathrm{SA}\!\left(
    \mathbf{H}^{\mathrm{fused}}
  \right)
  \label{eq:sa_final}
\end{equation}

\begin{equation}
  \mathbf{p} = \mathrm{Tanh}\!\left(
    \mathbf{W}_p\, \mathbf{H}^{\mathrm{out}}_{[0]}
  \right), \quad \mathbf{W}_p \in \mathbb{R}^{H \times H}
  \label{eq:pooler}
\end{equation}

\noindent trong đó $\mathbf{H}^{\mathrm{out}}_{[0]}$ là hàng ứng với
vị trí \texttt{[CLS]} (chỉ số 0).

\paragraph{Bước 4: Phân loại đa nhiệm.}
Từ biểu diễn câu $\mathbf{p} \in \mathbb{R}^H$, hai đầu ra được tính
song song:

\begin{align}
  \hat{\mathbf{y}}^{\mathrm{sent}} &= \mathbf{W}_s\, \mathbf{p},
    \quad \mathbf{W}_s \in \mathbb{R}^{3 \times H}
    \label{eq:head_sentiment}\\[4pt]
  \hat{\mathbf{y}}^{\mathrm{cat}} &= \mathbf{W}_c\, \mathbf{p},
    \quad \mathbf{W}_c \in \mathbb{R}^{6 \times H}
    \label{eq:head_category}
\end{align}

\noindent với $\hat{\mathbf{y}}^{\mathrm{sent}}$ là logit cho 3 lớp
cảm xúc (Positive / Negative / Neutral) và
$\hat{\mathbf{y}}^{\mathrm{cat}}$ là logit cho 6 nhóm khía cạnh
(SERVICE, FACILITY, AMENITY, EXPERIENCE, BRANDING, LOYALTY).

Hàm mất mát tổng hợp là tổng của hai hàm cross-entropy có trọng số
lớp:

\begin{equation}
  \mathcal{L} = \mathcal{L}_{\mathrm{CE}}^{\mathrm{sent}}
              + \mathcal{L}_{\mathrm{CE}}^{\mathrm{cat}}
  \label{eq:loss}
\end{equation}

Hình~\ref{fig:two_stream} minh họa toàn bộ kiến trúc hai luồng.

% =====================================================================
\section{Token Merging cho chuỗi 1D}
\label{sec:pp_tome}
% =====================================================================

\subsection{Tổng quan và vị trí tích hợp}
\label{subsec:tome_overview}

Token Merging (ToMe)~\cite{bolya2023tome} được tích hợp vào mô hình
\texttt{FastLcfBertMultiTask} dưới dạng module
\texttt{ToMeSequenceMerger}, áp dụng \textbf{sau} khi backbone BERT
hoàn thành quá trình mã hóa (post-BERT ToMe) và trước khi luồng cục bộ
được tính.

Module nhận đầu vào ba tensor:
$\mathbf{H} \in \mathbb{R}^{B \times L \times H}$ (trạng thái ẩn),
$\mathbf{V} \in \mathbb{R}^{B \times L}$ (véc-tơ chỉ thị khía cạnh)
và $\mathbf{M} \in \{0,1\}^{B \times L}$ (attention mask),
và trả về bộ $(\mathbf{H}', \mathbf{V}', \mathbf{M}')$ có độ dài
chuỗi $L' \leq L$.

Số vòng gộp (merge steps) được cấu hình qua
\texttt{num\_merge\_steps}$= 2$ (mặc định); mỗi vòng áp dụng một
trong ba chiến lược sau.

% ---------------------------------------------------------------------
\subsection{Phép gộp cơ bản và bảo toàn LCF}
\label{subsec:merge_op}
% ---------------------------------------------------------------------

Bất kể chiến lược nào, khi hai token $s$ (source) và $d$ (destination)
được chọn để gộp, phép gộp thực hiện hai thao tác:

\paragraph{Cập nhật biểu diễn bằng trung bình:}
\begin{equation}
  \mathbf{h}'_s = \frac{\mathbf{h}_s + \mathbf{h}_d}{2}
  \label{eq:merge_avg}
\end{equation}

\paragraph{Bảo toàn thông tin khía cạnh bằng max:}
\begin{equation}
  v'_s = \max(v_s,\; v_d)
  \label{eq:merge_lcf}
\end{equation}

Token $d$ bị loại khỏi chuỗi; token $s$ thay thế với giá trị đã được
cập nhật.  Quy tắc \eqref{eq:merge_lcf} đảm bảo rằng nếu một trong
hai token thuộc cụm từ khía cạnh (có $v = 1$), token hợp nhất kết quả
cũng mang nhãn khía cạnh, tránh mất tín hiệu vị trí cho cơ chế LCF
ở bước tiếp theo.

Tất cả ba chiến lược dùng \textbf{độ tương đồng cosine} trên biểu
diễn đã chuẩn hóa $L_2$ để đo mức độ giống nhau giữa hai token:

\begin{equation}
  \hat{\mathbf{h}}_i = \frac{\mathbf{h}_i}{\|\mathbf{h}_i\|_2 + \varepsilon},
  \quad \varepsilon = 10^{-6}
  \label{eq:l2norm}
\end{equation}

\begin{equation}
  \mathrm{sim}(i, j) = \hat{\mathbf{h}}_i \cdot \hat{\mathbf{h}}_j
  \label{eq:cosine}
\end{equation}

% ---------------------------------------------------------------------
\subsection{Chiến lược 1: Bipartite Matching}
\label{subsec:bipartite}
% ---------------------------------------------------------------------

Chiến lược này là chiến lược gốc của ToMe~\cite{bolya2023tome}, được
điều chỉnh cho chuỗi 1D và bổ sung bảo vệ token khía cạnh.

\paragraph{Xây dựng tập ứng viên.}
Từ $n$ token hợp lệ của câu (không tính padding), tập ứng viên gộp
$\mathcal{P}$ được xác định bằng cách loại trừ:
\begin{itemize}
  \item Token \texttt{[CLS]} ($i = 0$), bảo vệ bởi \texttt{protect\_cls};
  \item Token \texttt{[SEP]} cuối ($i = n-1$), bảo vệ bởi
    \texttt{protect\_sep};
  \item Các token khía cạnh ($v_i > 0.5$), bảo vệ bởi
    \texttt{protect\_aspect}.
\end{itemize}

\paragraph{Phân chia hai nhóm.}
Các phần tử của $\mathcal{P}$ được liệt kê theo thứ tự chỉ số tăng
dần thành dãy $p_0, p_1, \dots, p_{|\mathcal{P}|-1}$.  Dãy này được
chia xen kẽ thành hai nhóm:

\begin{align}
  \mathcal{A} &= \{p_0,\, p_2,\, p_4, \dots\} \quad \text{(chỉ số
    chẵn trong dãy)}\label{eq:group_A}\\
  \mathcal{B} &= \{p_1,\, p_3,\, p_5, \dots\} \quad \text{(chỉ số
    lẻ trong dãy)}\label{eq:group_B}
\end{align}

\paragraph{Tính ma trận độ tương đồng.}
Mỗi hàng ứng với một token trong $\mathcal{A}$, mỗi cột ứng với một
token trong $\mathcal{B}$:

\begin{equation}
  \mathbf{S} = \widehat{\mathbf{A}}\, \widehat{\mathbf{B}}^{\top},
  \quad \mathbf{S} \in \mathbb{R}^{|\mathcal{A}| \times |\mathcal{B}|}
  \label{eq:bipartite_score}
\end{equation}

\noindent trong đó $\widehat{\mathbf{A}}, \widehat{\mathbf{B}}$ là ma
trận biểu diễn chuẩn hóa $L_2$ của các token trong nhóm tương ứng.

\paragraph{Ghép cặp tham lam.}
Mỗi token $a_i \in \mathcal{A}$ chọn token $b_{j^*} \in \mathcal{B}$
có độ tương đồng lớn nhất:

\begin{equation}
  j^*(i) = \arg\max_{j} S_{ij}
  \label{eq:greedy_assign}
\end{equation}

Ràng buộc một-một phía $\mathcal{B}$: mỗi $b_j$ chỉ được gộp với một
$a_i$ duy nhất (cặp đầu tiên yêu cầu $b_j$ được giữ lại, các cặp sau
yêu cầu cùng $b_j$ bị bỏ qua).  Đây là phép ghép cặp hai phía
\textit{greedy} theo nghĩa mỗi bên $\mathcal{A}$ giải quyết độc lập
mà không có giai đoạn đấu giá toàn cục.

Sau bước này, mỗi cặp $(a_i, b_{j^*(i)})$ được gộp theo
\eqref{eq:merge_avg}--\eqref{eq:merge_lcf} với $s = a_i$, $d =
b_{j^*(i)}$.

\paragraph{Độ phức tạp.}
Với $|\mathcal{P}|$ token ứng viên và chiều ẩn $H$, chi phí tính ma
trận $\mathbf{S}$ là $\mathcal{O}(|\mathcal{P}|^2 / 4 \cdot H)$, thấp
hơn nhiều so với self-attention $\mathcal{O}(L^2 H)$ trên toàn chuỗi.

% ---------------------------------------------------------------------
\subsection{Chiến lược 2: Sequential Local Merging}
\label{subsec:sequential}
% ---------------------------------------------------------------------

Chiến lược này thực hiện gộp \textbf{tuần tự từ trái sang phải}, mỗi
token quyết định gộp dựa trên so sánh với hai lân cận gần nhất.

\paragraph{Xây dựng tập bảo vệ.}
Ngoài CLS, SEP và token khía cạnh, chiến lược này bổ sung bảo vệ
\textbf{ranh giới mid-SEP}: trong định dạng SPC
(xem~\eqref{eq:spc}), token phân cách giữa phần câu văn và phần
khía cạnh (tức \texttt{[SEP]} đầu tiên, không phải token cuối) không
được xóa.  Điều kiện nhận biết:

\begin{equation}
  \mathrm{midSep}(j) = \mathbf{1}\!\left[
    v_j < 0.5 \;\text{ và }\; v_{j+1} > 0.5
  \right]
  \label{eq:midsep}
\end{equation}

\paragraph{Vòng quét.}
Biến con trỏ $i$ duyệt từ $\texttt{protect\_left}$ đến $n -
\texttt{protect\_right} - 1$.  Tại mỗi bước $i$:

\begin{enumerate}
  \item Bỏ qua nếu $i$ đã bị xóa, là mid-SEP hoặc là token khía
    cạnh.
  \item Tìm \textbf{lân cận trái tích cực gần nhất}:
    $l = \max\{l' < i : l' \text{ chưa xóa và không là mid-SEP}\}$.
  \item Tìm \textbf{lân cận phải tích cực gần nhất không phải
    khía cạnh}:
    $r = \min\{r' > i : r' \text{ chưa xóa, không mid-SEP, không
    aspect}\}$.
  \item Tính độ tương đồng:
    \begin{align}
      \mathrm{sim}_{\mathrm{L}} &=
        \mathrm{sim}(i, l) \quad \text{(nếu tồn tại } l \text{)}
        \label{eq:sim_left}\\
      \mathrm{sim}_{\mathrm{R}} &=
        \mathrm{sim}(i, r) \quad \text{(nếu tồn tại } r \text{)}
        \label{eq:sim_right}
    \end{align}
  \item Ra quyết định gộp:
    \begin{equation}
      \begin{cases}
        \text{Token } i \text{ gộp vào trái}:
          & \mathbf{h}_l \leftarrow \tfrac{\mathbf{h}_l + \mathbf{h}_i}{2},\;
            v_l \leftarrow \max(v_l, v_i),\;
            \text{xóa } i \\
            & \quad \text{khi } \exists l \text{ và }
              \mathrm{sim}_{\mathrm{L}} > \mathrm{sim}_{\mathrm{R}}\\[6pt]
        \text{Token phải gộp vào } i:
          & \mathbf{h}_i \leftarrow \tfrac{\mathbf{h}_i + \mathbf{h}_r}{2},\;
            v_i \leftarrow \max(v_i, v_r),\;
            \text{xóa } r \\
            & \quad \text{khi } \exists r \text{ và }
              \mathrm{sim}_{\mathrm{R}} \geq \mathrm{sim}_{\mathrm{L}}
      \end{cases}
      \label{eq:seq_decision}
    \end{equation}
  \item Tăng $i \leftarrow i + 1$.
\end{enumerate}

Một lần duyệt đầy đủ ($i$ từ đầu đến cuối) được gọi là một
\textit{merge step}; mô hình thực hiện \texttt{num\_merge\_steps}$= 2$
lần duyệt liên tiếp.

\paragraph{Tính chất.}
Token có thể \textit{nhận} nhiều gộp từ nhiều hướng trong cùng một
lần duyệt (không có ràng buộc một-một phía đích).  Điều này khác với
Bipartite Matching và cho phép co chuỗi nhanh hơn khi có nhiều cụm
token lặp lại liên tiếp.

% ---------------------------------------------------------------------
\subsection{Chiến lược 3: Attention-Weighted Merging}
\label{subsec:attn_merge}
% ---------------------------------------------------------------------

Chiến lược này gộp ưu tiên \textbf{các token có trọng số chú ý thấp
nhất}, đồng thời bảo vệ các token mang thông tin quan trọng.

\paragraph{Xây dựng tập bảo vệ cứng.}
Ba nhóm token không bao giờ bị xóa:

\begin{itemize}
  \item \textbf{Biên cấu trúc}: CLS ($i = 0$) và SEP ($i = n-1$);
  \item \textbf{Token khía cạnh}: $v_i > 0.5$;
  \item \textbf{Token chú ý cao}: những token có trọng số chú ý
    $a_i$ nằm trong \textbf{phân vị thứ 75} trở lên:
    \begin{equation}
      \theta = \mathrm{quantile}(\mathbf{a},\; 0.75)
      \label{eq:attn_threshold}
    \end{equation}
    \begin{equation}
      \mathrm{highAttn}(i) = \mathbf{1}[a_i \geq \theta]
      \label{eq:high_attn}
    \end{equation}
\end{itemize}

Token $i$ bị bảo vệ khi thỏa ít nhất một trong các điều kiện trên.

\paragraph{Giới hạn số cặp gộp.}
\begin{equation}
  T_{\max} = \left\lfloor
    \frac{n - \texttt{protect\_left} - \texttt{protect\_right}}{2}
  \right\rfloor
  \label{eq:max_merges}
\end{equation}

\paragraph{Vòng lặp gộp lặp.}
Trong mỗi bước $t = 1, 2, \dots, T_{\max}$:

\begin{enumerate}
  \item \textbf{Tìm ứng viên bị xóa}: token chỉ số $i$ có trọng số
    chú ý nhỏ nhất trong số các token chưa bị xóa và không thuộc tập
    bảo vệ:
    \begin{equation}
      i = \arg\min_{\substack{t:\, \text{not removed}(t)\\
            \text{not protected}(t),\, v_t \leq 0.5}}
        a_t
      \label{eq:find_min_attn}
    \end{equation}
    Nếu không có ứng viên: dừng vòng lặp sớm.

  \item \textbf{Tìm token đích}: token $j \neq i$ chưa bị xóa, không
    phải khía cạnh, có độ tương đồng cosine với $i$ cao nhất:
    \begin{equation}
      j = \arg\max_{\substack{k \neq i:\, \text{not removed}(k)\\
            v_k \leq 0.5}}
        \hat{\mathbf{h}}_i \cdot \hat{\mathbf{h}}_k
      \label{eq:find_neighbor}
    \end{equation}
    Nếu giá trị tương đồng $< -1.5$ (không có lân cận hợp lệ): dừng.

  \item \textbf{Thực hiện gộp}: gộp $i$ vào $j$ theo
    \eqref{eq:merge_avg}--\eqref{eq:merge_lcf}, đánh dấu $i$ đã xóa.

  \item \textbf{Cập nhật metric}: chuẩn hóa lại toàn bộ ma trận biểu
    diễn trước vòng lặp kế tiếp.
\end{enumerate}

\paragraph{Sự khác biệt chính so với hai chiến lược còn lại.}
Bipartite Matching và Sequential Local chọn cặp dựa hoàn toàn vào độ
tương đồng cosine giữa biểu diễn.  Attention-Weighted Merging sử dụng
trọng số chú ý $\mathbf{a}$ như một \textit{tín hiệu tầm quan trọng}
độc lập để ưu tiên xóa các token mang ít thông tin cảm xúc nhất, đồng
thời bảo vệ không chỉ token khía cạnh mà cả những token có cơ chế chú
ý của BERT đã học được rằng chúng quan trọng.

% ---------------------------------------------------------------------
\subsection{Chế độ resize và compact}
\label{subsec:resize}
% ---------------------------------------------------------------------

Sau tất cả các vòng gộp, chuỗi có độ dài $L' \leq L$.  Module
\texttt{ToMeSequenceMerger} hỗ trợ hai chế độ đầu ra, điều khiển bởi
tham số \texttt{resize}:

\paragraph{Chế độ \texttt{resize=True} (nội suy về độ dài gốc).}
Chuỗi $\mathbf{H}' \in \mathbb{R}^{L' \times H}$ được nội suy tuyến
tính về đúng độ dài $L$ ban đầu:

\begin{equation}
  \mathbf{H}^{\mathrm{out}} =
    \mathrm{Interpolate}\!\left(\mathbf{H}',\; L,\;
    \text{mode=linear}\right)
    \in \mathbb{R}^{L \times H}
  \label{eq:resize}
\end{equation}

Chế độ này \textbf{không thay đổi hình dạng tensor} nên các lớp hạ
nguồn (SA, Linear) không cần điều chỉnh.  Nó cô lập tác động của
gộp token lên \textit{chất lượng biểu diễn} mà không đo tốc độ suy
luận.  Đây là chế độ được dùng trong toàn bộ các thực nghiệm chính của
luận văn.

\paragraph{Chế độ \texttt{resize=False} (compact).}
Chuỗi rút gọn $\mathbf{H}'$ được giữ nguyên; attention mask mới
$\mathbf{M}' \in \{0, 1\}^{B \times L'_{\max}}$ được tính và truyền
xuống các lớp self-attention.  Chế độ này mang lại \textbf{tăng tốc
thực sự} ($\mathcal{O}({L'}^2)$ thay cho $\mathcal{O}(L^2)$) nhưng
yêu cầu các lớp hạ nguồn hỗ trợ chuỗi có độ dài thay đổi.  Đây là
hướng tối ưu hóa được đề xuất cho các nghiên cứu tiếp theo.

\paragraph{Tóm tắt hai chế độ.}

\begin{table}[h]
  \centering
  \caption{So sánh hai chế độ đầu ra của \texttt{ToMeSequenceMerger}.}
  \label{tab:resize_mode}
  \begin{tabular}{lcc}
    \toprule
    \textbf{Thuộc tính} & \texttt{resize=True} & \texttt{resize=False} \\
    \midrule
    Độ dài đầu ra & $L$ (giữ nguyên) & $L' \leq L$ \\
    Hình dạng tensor thay đổi & Không & Có \\
    Tăng tốc suy luận thực sự & Không & Có \\
    Mục tiêu đánh giá & Chất lượng biểu diễn & Tốc độ + chất lượng \\
    Dùng trong thực nghiệm chính & \checkmark & \\
    \bottomrule
  \end{tabular}
\end{table}
 trong main.tex
% =====================================================================

% =====================================================================
\section{Cơ chế Local Context Focus}
\label{sec:pp_lcf}
% =====================================================================

\subsection{Định dạng đầu vào SPC}
\label{subsec:spc}

Mô hình \texttt{FastLcfBertMultiTask} nhận đầu vào theo định dạng
\textit{Single-Pair Classification} (SPC): cả câu đánh giá lẫn cụm từ
khía cạnh được nối thành một chuỗi duy nhất và đưa vào backbone BERT:

\begin{equation}
  \mathbf{X} = [\texttt{[CLS]},\; x_1, \dots, x_m,\;
                \texttt{[SEP]},\; a_1, \dots, a_k,\;
                \texttt{[SEP]}]
  \label{eq:spc}
\end{equation}

\noindent trong đó $x_1, \dots, x_m$ là các token của câu văn và
$a_1, \dots, a_k$ là các token của cụm từ khía cạnh.  Tổng độ dài
chuỗi là $L = m + k + 3$ (bao gồm hai token \texttt{[SEP]} và một
\texttt{[CLS]}).

Bên cạnh đó, một \textbf{véc-tơ chỉ thị khía cạnh} (aspect indicator)
$\mathbf{v} \in \{0, 1\}^{L}$ được xây dựng tương ứng:

\begin{equation}
  v_i = \begin{cases}
    1 & \text{nếu token thứ } i \text{ thuộc vị trí cụm khía cạnh
        (segment A, khớp với span gốc)}\\
    0 & \text{ngược lại}
  \end{cases}
  \label{eq:aspect_indicator}
\end{equation}

Véc-tơ $\mathbf{v}$ đóng vai trò đầu vào cho tất cả các cơ chế LCF và
được truyền xuyên suốt qua các module Token Merging (mục~\ref{sec:pp_tome}).

% ---------------------------------------------------------------------
\subsection{Khoảng cách ngữ nghĩa tương đối (SRD)}
\label{subsec:srd}
% ---------------------------------------------------------------------

Để định lượng mức độ ``gần'' của một token $i$ so với cụm từ khía
cạnh, luận văn sử dụng \textbf{Semantic Relative Distance} (SRD) theo
định nghĩa trong LCF-ATEPC~\cite{zeng2019lcf}:

\begin{equation}
  \mathrm{SRD}_i = \max\!\left(0,\;
    \left|i - c\right| - \left\lfloor \frac{l_a}{2} \right\rfloor
  \right)
  \label{eq:srd}
\end{equation}

\noindent với:

\begin{itemize}
  \item $c = \dfrac{a_{\mathrm{start}} + a_{\mathrm{end}}}{2}$: chỉ số
    trung tâm của cụm từ khía cạnh trong chuỗi token;
  \item $l_a = a_{\mathrm{end}} - a_{\mathrm{start}} + 1$: độ dài (số
    token) của cụm từ khía cạnh;
  \item $a_{\mathrm{start}}, a_{\mathrm{end}}$: chỉ số token đầu và
    cuối của cụm khía cạnh trong chuỗi đầy đủ.
\end{itemize}

Khi token $i$ nằm trong vùng trực tiếp bao phủ bởi cụm từ khía cạnh,
$\mathrm{SRD}_i = 0$.  Khi token nằm ngoài vùng đó, $\mathrm{SRD}_i$
tăng tuyến tính theo khoảng cách từ biên cụm khía cạnh đến token $i$.

Ngưỡng bán kính $\alpha$ (tham số \texttt{srd\_threshold}, mặc định
$\alpha = 5$) xác định vùng ngữ cảnh cục bộ: mọi token có
$\mathrm{SRD}_i \leq \alpha$ được coi là \textit{nằm trong ngữ cảnh
cục bộ} và nhận trọng số tối đa bằng 1.

% ---------------------------------------------------------------------
\subsection{Context Dynamic Weighting (CDW)}
\label{subsec:cdw}
% ---------------------------------------------------------------------

CDW gán cho mỗi token một \textbf{trọng số liên tục} $w_i \in [0, 1]$
phản ánh mức độ liên quan đến khía cạnh đang xét:

\begin{equation}
  w_i^{\mathrm{CDW}} = \begin{cases}
    1 & \text{nếu } \mathrm{SRD}_i \leq \alpha \\[6pt]
    \dfrac{L - (\mathrm{SRD}_i - \alpha)}{L}
      & \text{nếu } \mathrm{SRD}_i > \alpha
  \end{cases}
  \label{eq:cdw}
\end{equation}

\noindent sau đó được kẹp vào đoạn $[0, 1]$:

\begin{equation}
  w_i^{\mathrm{CDW}} \leftarrow \mathrm{clamp}\!\left(
    w_i^{\mathrm{CDW}},\; 0,\; 1
  \right)
  \label{eq:cdw_clamp}
\end{equation}

Vùng trong bán kính $\alpha$ nhận trọng số 1 (không suy giảm); càng
ra xa khía cạnh, trọng số giảm tuyến tính và bằng 0 khi
$\mathrm{SRD}_i = \alpha + L$.  Trong trường hợp không có token khía
cạnh nào được phát hiện trong câu (véc-tơ $\mathbf{v}$ toàn 0), CDW
trả về véc-tơ trọng số hằng số bằng 1 cho toàn bộ chuỗi.

Véc-tơ trọng số CDW $\mathbf{w}^{\mathrm{CDW}} \in \mathbb{R}^L$ được
nhân phần tử với ma trận trạng thái ẩn của BERT để tạo ra
\textbf{luồng cục bộ} (local stream):

\begin{equation}
  \mathbf{H}^{\mathrm{local}} = \mathbf{H} \odot
    \mathbf{w}^{\mathrm{CDW}} \mathbf{1}_{H}^{\top}
  \label{eq:local_cdw}
\end{equation}

\noindent trong đó $\mathbf{H} \in \mathbb{R}^{L \times H}$ là đầu ra
của backbone BERT, $\odot$ là tích Hadamard theo chiều chuỗi (mỗi
hàng $\mathbf{H}_i$ được nhân với vô hướng $w_i$), và
$\mathbf{1}_{H}^{\top}$ ký hiệu phép broadcast sang chiều ẩn $H$.

% ---------------------------------------------------------------------
\subsection{Context Dynamic Masking (CDM)}
\label{subsec:cdm}
% ---------------------------------------------------------------------

CDM là biến thể nhị phân của CDW: thay vì giảm trọng số liên tục, nó
\textbf{triệt tiêu hoàn toàn} những token nằm ngoài ngưỡng $\alpha$:

\begin{equation}
  w_i^{\mathrm{CDM}} = \begin{cases}
    1 & \text{nếu } \mathrm{SRD}_i \leq \alpha \\[4pt]
    0 & \text{nếu } \mathrm{SRD}_i > \alpha
  \end{cases}
  \label{eq:cdm}
\end{equation}

So với CDW, mặt nạ CDM tạo ra một ranh giới cứng quanh vùng ngữ cảnh
cục bộ. Luồng cục bộ khi dùng CDM là:

\begin{equation}
  \mathbf{H}^{\mathrm{local}} = \mathbf{H} \odot
    \mathbf{w}^{\mathrm{CDM}} \mathbf{1}_{H}^{\top}
  \label{eq:local_cdm}
\end{equation}

Khi không có token khía cạnh, CDM cũng trả về véc-tơ trọng số toàn 1
(tương tự CDW), đảm bảo mô hình không suy sụp hoàn toàn trên các mẫu
thiếu annotation.

% ---------------------------------------------------------------------
\subsection{Kiến trúc hai luồng và phân loại đa nhiệm}
\label{subsec:two_stream}
% ---------------------------------------------------------------------

Sau khi tính luồng cục bộ bằng CDW hoặc CDM, mô hình thực hiện
\textbf{thiết kế hai luồng} (two-stream) kết hợp thông tin cục bộ
và toàn cục:

\paragraph{Bước 1: Self-attention cục bộ.}
Luồng cục bộ được đưa qua một lớp self-attention bổ sung
$\mathrm{SA}_{\mathrm{LCF}}$ với kết nối dư (residual) và
chuẩn hóa lớp (LayerNorm):

\begin{equation}
  \widetilde{\mathbf{H}}^{\mathrm{local}} =
    \mathrm{LayerNorm}\!\left(
      \mathbf{H}^{\mathrm{local}} +
      \mathrm{Dropout}\!\left(
        \mathrm{MHA}\!\left(
          \mathbf{H}^{\mathrm{local}},\,
          \mathbf{H}^{\mathrm{local}},\,
          \mathbf{H}^{\mathrm{local}}
        \right)
      \right)
    \right)
  \label{eq:sa_lcf}
\end{equation}

\noindent trong đó $\mathrm{MHA}(\cdot)$ là phép multi-head attention
với 8 head.

\paragraph{Bước 2: Hợp nhất hai luồng.}
Luồng cục bộ và luồng toàn cục được ghép nối theo chiều ẩn rồi chiếu
xuống $H$ chiều:

\begin{equation}
  \mathbf{H}^{\mathrm{fused}} = \mathrm{Dropout}\!\left(
    \mathbf{W}_2 \left[
      \widetilde{\mathbf{H}}^{\mathrm{local}};
      \mathbf{H}
    \right]
  \right), \quad \mathbf{W}_2 \in \mathbb{R}^{H \times 2H}
  \label{eq:fusion}
\end{equation}

\paragraph{Bước 3: Self-attention tổng hợp và pooling.}
Biểu diễn tổng hợp tiếp tục qua một lớp self-attention $\mathrm{SA}$
nữa trước khi lấy biểu diễn token \texttt{[CLS]}:

\begin{equation}
  \mathbf{H}^{\mathrm{out}} = \mathrm{SA}\!\left(
    \mathbf{H}^{\mathrm{fused}}
  \right)
  \label{eq:sa_final}
\end{equation}

\begin{equation}
  \mathbf{p} = \mathrm{Tanh}\!\left(
    \mathbf{W}_p\, \mathbf{H}^{\mathrm{out}}_{[0]}
  \right), \quad \mathbf{W}_p \in \mathbb{R}^{H \times H}
  \label{eq:pooler}
\end{equation}

\noindent trong đó $\mathbf{H}^{\mathrm{out}}_{[0]}$ là hàng ứng với
vị trí \texttt{[CLS]} (chỉ số 0).

\paragraph{Bước 4: Phân loại đa nhiệm.}
Từ biểu diễn câu $\mathbf{p} \in \mathbb{R}^H$, hai đầu ra được tính
song song:

\begin{align}
  \hat{\mathbf{y}}^{\mathrm{sent}} &= \mathbf{W}_s\, \mathbf{p},
    \quad \mathbf{W}_s \in \mathbb{R}^{3 \times H}
    \label{eq:head_sentiment}\\[4pt]
  \hat{\mathbf{y}}^{\mathrm{cat}} &= \mathbf{W}_c\, \mathbf{p},
    \quad \mathbf{W}_c \in \mathbb{R}^{6 \times H}
    \label{eq:head_category}
\end{align}

\noindent với $\hat{\mathbf{y}}^{\mathrm{sent}}$ là logit cho 3 lớp
cảm xúc (Positive / Negative / Neutral) và
$\hat{\mathbf{y}}^{\mathrm{cat}}$ là logit cho 6 nhóm khía cạnh
(SERVICE, FACILITY, AMENITY, EXPERIENCE, BRANDING, LOYALTY).

Hàm mất mát tổng hợp là tổng của hai hàm cross-entropy có trọng số
lớp:

\begin{equation}
  \mathcal{L} = \mathcal{L}_{\mathrm{CE}}^{\mathrm{sent}}
              + \mathcal{L}_{\mathrm{CE}}^{\mathrm{cat}}
  \label{eq:loss}
\end{equation}

Hình~\ref{fig:two_stream} minh họa toàn bộ kiến trúc hai luồng.

% =====================================================================
\section{Token Merging cho chuỗi 1D}
\label{sec:pp_tome}
% =====================================================================

\subsection{Tổng quan và vị trí tích hợp}
\label{subsec:tome_overview}

Token Merging (ToMe)~\cite{bolya2023tome} được tích hợp vào mô hình
\texttt{FastLcfBertMultiTask} dưới dạng module
\texttt{ToMeSequenceMerger}, áp dụng \textbf{sau} khi backbone BERT
hoàn thành quá trình mã hóa (post-BERT ToMe) và trước khi luồng cục bộ
được tính.

Module nhận đầu vào ba tensor:
$\mathbf{H} \in \mathbb{R}^{B \times L \times H}$ (trạng thái ẩn),
$\mathbf{V} \in \mathbb{R}^{B \times L}$ (véc-tơ chỉ thị khía cạnh)
và $\mathbf{M} \in \{0,1\}^{B \times L}$ (attention mask),
và trả về bộ $(\mathbf{H}', \mathbf{V}', \mathbf{M}')$ có độ dài
chuỗi $L' \leq L$.

Số vòng gộp (merge steps) được cấu hình qua
\texttt{num\_merge\_steps}$= 2$ (mặc định); mỗi vòng áp dụng một
trong ba chiến lược sau.

% ---------------------------------------------------------------------
\subsection{Phép gộp cơ bản và bảo toàn LCF}
\label{subsec:merge_op}
% ---------------------------------------------------------------------

Bất kể chiến lược nào, khi hai token $s$ (source) và $d$ (destination)
được chọn để gộp, phép gộp thực hiện hai thao tác:

\paragraph{Cập nhật biểu diễn bằng trung bình:}
\begin{equation}
  \mathbf{h}'_s = \frac{\mathbf{h}_s + \mathbf{h}_d}{2}
  \label{eq:merge_avg}
\end{equation}

\paragraph{Bảo toàn thông tin khía cạnh bằng max:}
\begin{equation}
  v'_s = \max(v_s,\; v_d)
  \label{eq:merge_lcf}
\end{equation}

Token $d$ bị loại khỏi chuỗi; token $s$ thay thế với giá trị đã được
cập nhật.  Quy tắc \eqref{eq:merge_lcf} đảm bảo rằng nếu một trong
hai token thuộc cụm từ khía cạnh (có $v = 1$), token hợp nhất kết quả
cũng mang nhãn khía cạnh, tránh mất tín hiệu vị trí cho cơ chế LCF
ở bước tiếp theo.

Tất cả ba chiến lược dùng \textbf{độ tương đồng cosine} trên biểu
diễn đã chuẩn hóa $L_2$ để đo mức độ giống nhau giữa hai token:

\begin{equation}
  \hat{\mathbf{h}}_i = \frac{\mathbf{h}_i}{\|\mathbf{h}_i\|_2 + \varepsilon},
  \quad \varepsilon = 10^{-6}
  \label{eq:l2norm}
\end{equation}

\begin{equation}
  \mathrm{sim}(i, j) = \hat{\mathbf{h}}_i \cdot \hat{\mathbf{h}}_j
  \label{eq:cosine}
\end{equation}

% ---------------------------------------------------------------------
\subsection{Chiến lược 1: Bipartite Matching}
\label{subsec:bipartite}
% ---------------------------------------------------------------------

Chiến lược này là chiến lược gốc của ToMe~\cite{bolya2023tome}, được
điều chỉnh cho chuỗi 1D và bổ sung bảo vệ token khía cạnh.

\paragraph{Xây dựng tập ứng viên.}
Từ $n$ token hợp lệ của câu (không tính padding), tập ứng viên gộp
$\mathcal{P}$ được xác định bằng cách loại trừ:
\begin{itemize}
  \item Token \texttt{[CLS]} ($i = 0$), bảo vệ bởi \texttt{protect\_cls};
  \item Token \texttt{[SEP]} cuối ($i = n-1$), bảo vệ bởi
    \texttt{protect\_sep};
  \item Các token khía cạnh ($v_i > 0.5$), bảo vệ bởi
    \texttt{protect\_aspect}.
\end{itemize}

\paragraph{Phân chia hai nhóm.}
Các phần tử của $\mathcal{P}$ được liệt kê theo thứ tự chỉ số tăng
dần thành dãy $p_0, p_1, \dots, p_{|\mathcal{P}|-1}$.  Dãy này được
chia xen kẽ thành hai nhóm:

\begin{align}
  \mathcal{A} &= \{p_0,\, p_2,\, p_4, \dots\} \quad \text{(chỉ số
    chẵn trong dãy)}\label{eq:group_A}\\
  \mathcal{B} &= \{p_1,\, p_3,\, p_5, \dots\} \quad \text{(chỉ số
    lẻ trong dãy)}\label{eq:group_B}
\end{align}

\paragraph{Tính ma trận độ tương đồng.}
Mỗi hàng ứng với một token trong $\mathcal{A}$, mỗi cột ứng với một
token trong $\mathcal{B}$:

\begin{equation}
  \mathbf{S} = \widehat{\mathbf{A}}\, \widehat{\mathbf{B}}^{\top},
  \quad \mathbf{S} \in \mathbb{R}^{|\mathcal{A}| \times |\mathcal{B}|}
  \label{eq:bipartite_score}
\end{equation}

\noindent trong đó $\widehat{\mathbf{A}}, \widehat{\mathbf{B}}$ là ma
trận biểu diễn chuẩn hóa $L_2$ của các token trong nhóm tương ứng.

\paragraph{Ghép cặp tham lam.}
Mỗi token $a_i \in \mathcal{A}$ chọn token $b_{j^*} \in \mathcal{B}$
có độ tương đồng lớn nhất:

\begin{equation}
  j^*(i) = \arg\max_{j} S_{ij}
  \label{eq:greedy_assign}
\end{equation}

Ràng buộc một-một phía $\mathcal{B}$: mỗi $b_j$ chỉ được gộp với một
$a_i$ duy nhất (cặp đầu tiên yêu cầu $b_j$ được giữ lại, các cặp sau
yêu cầu cùng $b_j$ bị bỏ qua).  Đây là phép ghép cặp hai phía
\textit{greedy} theo nghĩa mỗi bên $\mathcal{A}$ giải quyết độc lập
mà không có giai đoạn đấu giá toàn cục.

Sau bước này, mỗi cặp $(a_i, b_{j^*(i)})$ được gộp theo
\eqref{eq:merge_avg}--\eqref{eq:merge_lcf} với $s = a_i$, $d =
b_{j^*(i)}$.

\paragraph{Độ phức tạp.}
Với $|\mathcal{P}|$ token ứng viên và chiều ẩn $H$, chi phí tính ma
trận $\mathbf{S}$ là $\mathcal{O}(|\mathcal{P}|^2 / 4 \cdot H)$, thấp
hơn nhiều so với self-attention $\mathcal{O}(L^2 H)$ trên toàn chuỗi.

% ---------------------------------------------------------------------
\subsection{Chiến lược 2: Sequential Local Merging}
\label{subsec:sequential}
% ---------------------------------------------------------------------

Chiến lược này thực hiện gộp \textbf{tuần tự từ trái sang phải}, mỗi
token quyết định gộp dựa trên so sánh với hai lân cận gần nhất.

\paragraph{Xây dựng tập bảo vệ.}
Ngoài CLS, SEP và token khía cạnh, chiến lược này bổ sung bảo vệ
\textbf{ranh giới mid-SEP}: trong định dạng SPC
(xem~\eqref{eq:spc}), token phân cách giữa phần câu văn và phần
khía cạnh (tức \texttt{[SEP]} đầu tiên, không phải token cuối) không
được xóa.  Điều kiện nhận biết:

\begin{equation}
  \mathrm{midSep}(j) = \mathbf{1}\!\left[
    v_j < 0.5 \;\text{ và }\; v_{j+1} > 0.5
  \right]
  \label{eq:midsep}
\end{equation}

\paragraph{Vòng quét.}
Biến con trỏ $i$ duyệt từ $\texttt{protect\_left}$ đến $n -
\texttt{protect\_right} - 1$.  Tại mỗi bước $i$:

\begin{enumerate}
  \item Bỏ qua nếu $i$ đã bị xóa, là mid-SEP hoặc là token khía
    cạnh.
  \item Tìm \textbf{lân cận trái tích cực gần nhất}:
    $l = \max\{l' < i : l' \text{ chưa xóa và không là mid-SEP}\}$.
  \item Tìm \textbf{lân cận phải tích cực gần nhất không phải
    khía cạnh}:
    $r = \min\{r' > i : r' \text{ chưa xóa, không mid-SEP, không
    aspect}\}$.
  \item Tính độ tương đồng:
    \begin{align}
      \mathrm{sim}_{\mathrm{L}} &=
        \mathrm{sim}(i, l) \quad \text{(nếu tồn tại } l \text{)}
        \label{eq:sim_left}\\
      \mathrm{sim}_{\mathrm{R}} &=
        \mathrm{sim}(i, r) \quad \text{(nếu tồn tại } r \text{)}
        \label{eq:sim_right}
    \end{align}
  \item Ra quyết định gộp:
    \begin{equation}
      \begin{cases}
        \text{Token } i \text{ gộp vào trái}:
          & \mathbf{h}_l \leftarrow \tfrac{\mathbf{h}_l + \mathbf{h}_i}{2},\;
            v_l \leftarrow \max(v_l, v_i),\;
            \text{xóa } i \\
            & \quad \text{khi } \exists l \text{ và }
              \mathrm{sim}_{\mathrm{L}} > \mathrm{sim}_{\mathrm{R}}\\[6pt]
        \text{Token phải gộp vào } i:
          & \mathbf{h}_i \leftarrow \tfrac{\mathbf{h}_i + \mathbf{h}_r}{2},\;
            v_i \leftarrow \max(v_i, v_r),\;
            \text{xóa } r \\
            & \quad \text{khi } \exists r \text{ và }
              \mathrm{sim}_{\mathrm{R}} \geq \mathrm{sim}_{\mathrm{L}}
      \end{cases}
      \label{eq:seq_decision}
    \end{equation}
  \item Tăng $i \leftarrow i + 1$.
\end{enumerate}

Một lần duyệt đầy đủ ($i$ từ đầu đến cuối) được gọi là một
\textit{merge step}; mô hình thực hiện \texttt{num\_merge\_steps}$= 2$
lần duyệt liên tiếp.

\paragraph{Tính chất.}
Token có thể \textit{nhận} nhiều gộp từ nhiều hướng trong cùng một
lần duyệt (không có ràng buộc một-một phía đích).  Điều này khác với
Bipartite Matching và cho phép co chuỗi nhanh hơn khi có nhiều cụm
token lặp lại liên tiếp.

% ---------------------------------------------------------------------
\subsection{Chiến lược 3: Attention-Weighted Merging}
\label{subsec:attn_merge}
% ---------------------------------------------------------------------

Chiến lược này gộp ưu tiên \textbf{các token có trọng số chú ý thấp
nhất}, đồng thời bảo vệ các token mang thông tin quan trọng.

\paragraph{Xây dựng tập bảo vệ cứng.}
Ba nhóm token không bao giờ bị xóa:

\begin{itemize}
  \item \textbf{Biên cấu trúc}: CLS ($i = 0$) và SEP ($i = n-1$);
  \item \textbf{Token khía cạnh}: $v_i > 0.5$;
  \item \textbf{Token chú ý cao}: những token có trọng số chú ý
    $a_i$ nằm trong \textbf{phân vị thứ 75} trở lên:
    \begin{equation}
      \theta = \mathrm{quantile}(\mathbf{a},\; 0.75)
      \label{eq:attn_threshold}
    \end{equation}
    \begin{equation}
      \mathrm{highAttn}(i) = \mathbf{1}[a_i \geq \theta]
      \label{eq:high_attn}
    \end{equation}
\end{itemize}

Token $i$ bị bảo vệ khi thỏa ít nhất một trong các điều kiện trên.

\paragraph{Giới hạn số cặp gộp.}
\begin{equation}
  T_{\max} = \left\lfloor
    \frac{n - \texttt{protect\_left} - \texttt{protect\_right}}{2}
  \right\rfloor
  \label{eq:max_merges}
\end{equation}

\paragraph{Vòng lặp gộp lặp.}
Trong mỗi bước $t = 1, 2, \dots, T_{\max}$:

\begin{enumerate}
  \item \textbf{Tìm ứng viên bị xóa}: token chỉ số $i$ có trọng số
    chú ý nhỏ nhất trong số các token chưa bị xóa và không thuộc tập
    bảo vệ:
    \begin{equation}
      i = \arg\min_{\substack{t:\, \text{not removed}(t)\\
            \text{not protected}(t),\, v_t \leq 0.5}}
        a_t
      \label{eq:find_min_attn}
    \end{equation}
    Nếu không có ứng viên: dừng vòng lặp sớm.

  \item \textbf{Tìm token đích}: token $j \neq i$ chưa bị xóa, không
    phải khía cạnh, có độ tương đồng cosine với $i$ cao nhất:
    \begin{equation}
      j = \arg\max_{\substack{k \neq i:\, \text{not removed}(k)\\
            v_k \leq 0.5}}
        \hat{\mathbf{h}}_i \cdot \hat{\mathbf{h}}_k
      \label{eq:find_neighbor}
    \end{equation}
    Nếu giá trị tương đồng $< -1.5$ (không có lân cận hợp lệ): dừng.

  \item \textbf{Thực hiện gộp}: gộp $i$ vào $j$ theo
    \eqref{eq:merge_avg}--\eqref{eq:merge_lcf}, đánh dấu $i$ đã xóa.

  \item \textbf{Cập nhật metric}: chuẩn hóa lại toàn bộ ma trận biểu
    diễn trước vòng lặp kế tiếp.
\end{enumerate}

\paragraph{Sự khác biệt chính so với hai chiến lược còn lại.}
Bipartite Matching và Sequential Local chọn cặp dựa hoàn toàn vào độ
tương đồng cosine giữa biểu diễn.  Attention-Weighted Merging sử dụng
trọng số chú ý $\mathbf{a}$ như một \textit{tín hiệu tầm quan trọng}
độc lập để ưu tiên xóa các token mang ít thông tin cảm xúc nhất, đồng
thời bảo vệ không chỉ token khía cạnh mà cả những token có cơ chế chú
ý của BERT đã học được rằng chúng quan trọng.

% ---------------------------------------------------------------------
\subsection{Chế độ resize và compact}
\label{subsec:resize}
% ---------------------------------------------------------------------

Sau tất cả các vòng gộp, chuỗi có độ dài $L' \leq L$.  Module
\texttt{ToMeSequenceMerger} hỗ trợ hai chế độ đầu ra, điều khiển bởi
tham số \texttt{resize}:

\paragraph{Chế độ \texttt{resize=True} (nội suy về độ dài gốc).}
Chuỗi $\mathbf{H}' \in \mathbb{R}^{L' \times H}$ được nội suy tuyến
tính về đúng độ dài $L$ ban đầu:

\begin{equation}
  \mathbf{H}^{\mathrm{out}} =
    \mathrm{Interpolate}\!\left(\mathbf{H}',\; L,\;
    \text{mode=linear}\right)
    \in \mathbb{R}^{L \times H}
  \label{eq:resize}
\end{equation}

Chế độ này \textbf{không thay đổi hình dạng tensor} nên các lớp hạ
nguồn (SA, Linear) không cần điều chỉnh.  Nó cô lập tác động của
gộp token lên \textit{chất lượng biểu diễn} mà không đo tốc độ suy
luận.  Đây là chế độ được dùng trong toàn bộ các thực nghiệm chính của
luận văn.

\paragraph{Chế độ \texttt{resize=False} (compact).}
Chuỗi rút gọn $\mathbf{H}'$ được giữ nguyên; attention mask mới
$\mathbf{M}' \in \{0, 1\}^{B \times L'_{\max}}$ được tính và truyền
xuống các lớp self-attention.  Chế độ này mang lại \textbf{tăng tốc
thực sự} ($\mathcal{O}({L'}^2)$ thay cho $\mathcal{O}(L^2)$) nhưng
yêu cầu các lớp hạ nguồn hỗ trợ chuỗi có độ dài thay đổi.  Đây là
hướng tối ưu hóa được đề xuất cho các nghiên cứu tiếp theo.

\paragraph{Tóm tắt hai chế độ.}

\begin{table}[h]
  \centering
  \caption{So sánh hai chế độ đầu ra của \texttt{ToMeSequenceMerger}.}
  \label{tab:resize_mode}
  \begin{tabular}{lcc}
    \toprule
    \textbf{Thuộc tính} & \texttt{resize=True} & \texttt{resize=False} \\
    \midrule
    Độ dài đầu ra & $L$ (giữ nguyên) & $L' \leq L$ \\
    Hình dạng tensor thay đổi & Không & Có \\
    Tăng tốc suy luận thực sự & Không & Có \\
    Mục tiêu đánh giá & Chất lượng biểu diễn & Tốc độ + chất lượng \\
    Dùng trong thực nghiệm chính & \checkmark & \\
    \bottomrule
  \end{tabular}
\end{table}
 trong main.tex
% =====================================================================

% =====================================================================
\section{Cơ chế Local Context Focus}
\label{sec:pp_lcf}
% =====================================================================

\subsection{Định dạng đầu vào SPC}
\label{subsec:spc}

Mô hình \texttt{FastLcfBertMultiTask} nhận đầu vào theo định dạng
\textit{Single-Pair Classification} (SPC): cả câu đánh giá lẫn cụm từ
khía cạnh được nối thành một chuỗi duy nhất và đưa vào backbone BERT:

\begin{equation}
  \mathbf{X} = [\texttt{[CLS]},\; x_1, \dots, x_m,\;
                \texttt{[SEP]},\; a_1, \dots, a_k,\;
                \texttt{[SEP]}]
  \label{eq:spc}
\end{equation}

\noindent trong đó $x_1, \dots, x_m$ là các token của câu văn và
$a_1, \dots, a_k$ là các token của cụm từ khía cạnh.  Tổng độ dài
chuỗi là $L = m + k + 3$ (bao gồm hai token \texttt{[SEP]} và một
\texttt{[CLS]}).

Bên cạnh đó, một \textbf{véc-tơ chỉ thị khía cạnh} (aspect indicator)
$\mathbf{v} \in \{0, 1\}^{L}$ được xây dựng tương ứng:

\begin{equation}
  v_i = \begin{cases}
    1 & \text{nếu token thứ } i \text{ thuộc vị trí cụm khía cạnh
        (segment A, khớp với span gốc)}\\
    0 & \text{ngược lại}
  \end{cases}
  \label{eq:aspect_indicator}
\end{equation}

Véc-tơ $\mathbf{v}$ đóng vai trò đầu vào cho tất cả các cơ chế LCF và
được truyền xuyên suốt qua các module Token Merging (mục~\ref{sec:pp_tome}).

% ---------------------------------------------------------------------
\subsection{Khoảng cách ngữ nghĩa tương đối (SRD)}
\label{subsec:srd}
% ---------------------------------------------------------------------

Để định lượng mức độ ``gần'' của một token $i$ so với cụm từ khía
cạnh, luận văn sử dụng \textbf{Semantic Relative Distance} (SRD) theo
định nghĩa trong LCF-ATEPC~\cite{zeng2019lcf}:

\begin{equation}
  \mathrm{SRD}_i = \max\!\left(0,\;
    \left|i - c\right| - \left\lfloor \frac{l_a}{2} \right\rfloor
  \right)
  \label{eq:srd}
\end{equation}

\noindent với:

\begin{itemize}
  \item $c = \dfrac{a_{\mathrm{start}} + a_{\mathrm{end}}}{2}$: chỉ số
    trung tâm của cụm từ khía cạnh trong chuỗi token;
  \item $l_a = a_{\mathrm{end}} - a_{\mathrm{start}} + 1$: độ dài (số
    token) của cụm từ khía cạnh;
  \item $a_{\mathrm{start}}, a_{\mathrm{end}}$: chỉ số token đầu và
    cuối của cụm khía cạnh trong chuỗi đầy đủ.
\end{itemize}

Khi token $i$ nằm trong vùng trực tiếp bao phủ bởi cụm từ khía cạnh,
$\mathrm{SRD}_i = 0$.  Khi token nằm ngoài vùng đó, $\mathrm{SRD}_i$
tăng tuyến tính theo khoảng cách từ biên cụm khía cạnh đến token $i$.

Ngưỡng bán kính $\alpha$ (tham số \texttt{srd\_threshold}, mặc định
$\alpha = 5$) xác định vùng ngữ cảnh cục bộ: mọi token có
$\mathrm{SRD}_i \leq \alpha$ được coi là \textit{nằm trong ngữ cảnh
cục bộ} và nhận trọng số tối đa bằng 1.

% ---------------------------------------------------------------------
\subsection{Context Dynamic Weighting (CDW)}
\label{subsec:cdw}
% ---------------------------------------------------------------------

CDW gán cho mỗi token một \textbf{trọng số liên tục} $w_i \in [0, 1]$
phản ánh mức độ liên quan đến khía cạnh đang xét:

\begin{equation}
  w_i^{\mathrm{CDW}} = \begin{cases}
    1 & \text{nếu } \mathrm{SRD}_i \leq \alpha \\[6pt]
    \dfrac{L - (\mathrm{SRD}_i - \alpha)}{L}
      & \text{nếu } \mathrm{SRD}_i > \alpha
  \end{cases}
  \label{eq:cdw}
\end{equation}

\noindent sau đó được kẹp vào đoạn $[0, 1]$:

\begin{equation}
  w_i^{\mathrm{CDW}} \leftarrow \mathrm{clamp}\!\left(
    w_i^{\mathrm{CDW}},\; 0,\; 1
  \right)
  \label{eq:cdw_clamp}
\end{equation}

Vùng trong bán kính $\alpha$ nhận trọng số 1 (không suy giảm); càng
ra xa khía cạnh, trọng số giảm tuyến tính và bằng 0 khi
$\mathrm{SRD}_i = \alpha + L$.  Trong trường hợp không có token khía
cạnh nào được phát hiện trong câu (véc-tơ $\mathbf{v}$ toàn 0), CDW
trả về véc-tơ trọng số hằng số bằng 1 cho toàn bộ chuỗi.

Véc-tơ trọng số CDW $\mathbf{w}^{\mathrm{CDW}} \in \mathbb{R}^L$ được
nhân phần tử với ma trận trạng thái ẩn của BERT để tạo ra
\textbf{luồng cục bộ} (local stream):

\begin{equation}
  \mathbf{H}^{\mathrm{local}} = \mathbf{H} \odot
    \mathbf{w}^{\mathrm{CDW}} \mathbf{1}_{H}^{\top}
  \label{eq:local_cdw}
\end{equation}

\noindent trong đó $\mathbf{H} \in \mathbb{R}^{L \times H}$ là đầu ra
của backbone BERT, $\odot$ là tích Hadamard theo chiều chuỗi (mỗi
hàng $\mathbf{H}_i$ được nhân với vô hướng $w_i$), và
$\mathbf{1}_{H}^{\top}$ ký hiệu phép broadcast sang chiều ẩn $H$.

% ---------------------------------------------------------------------
\subsection{Context Dynamic Masking (CDM)}
\label{subsec:cdm}
% ---------------------------------------------------------------------

CDM là biến thể nhị phân của CDW: thay vì giảm trọng số liên tục, nó
\textbf{triệt tiêu hoàn toàn} những token nằm ngoài ngưỡng $\alpha$:

\begin{equation}
  w_i^{\mathrm{CDM}} = \begin{cases}
    1 & \text{nếu } \mathrm{SRD}_i \leq \alpha \\[4pt]
    0 & \text{nếu } \mathrm{SRD}_i > \alpha
  \end{cases}
  \label{eq:cdm}
\end{equation}

So với CDW, mặt nạ CDM tạo ra một ranh giới cứng quanh vùng ngữ cảnh
cục bộ. Luồng cục bộ khi dùng CDM là:

\begin{equation}
  \mathbf{H}^{\mathrm{local}} = \mathbf{H} \odot
    \mathbf{w}^{\mathrm{CDM}} \mathbf{1}_{H}^{\top}
  \label{eq:local_cdm}
\end{equation}

Khi không có token khía cạnh, CDM cũng trả về véc-tơ trọng số toàn 1
(tương tự CDW), đảm bảo mô hình không suy sụp hoàn toàn trên các mẫu
thiếu annotation.

% ---------------------------------------------------------------------
\subsection{Kiến trúc hai luồng và phân loại đa nhiệm}
\label{subsec:two_stream}
% ---------------------------------------------------------------------

Sau khi tính luồng cục bộ bằng CDW hoặc CDM, mô hình thực hiện
\textbf{thiết kế hai luồng} (two-stream) kết hợp thông tin cục bộ
và toàn cục:

\paragraph{Bước 1: Self-attention cục bộ.}
Luồng cục bộ được đưa qua một lớp self-attention bổ sung
$\mathrm{SA}_{\mathrm{LCF}}$ với kết nối dư (residual) và
chuẩn hóa lớp (LayerNorm):

\begin{equation}
  \widetilde{\mathbf{H}}^{\mathrm{local}} =
    \mathrm{LayerNorm}\!\left(
      \mathbf{H}^{\mathrm{local}} +
      \mathrm{Dropout}\!\left(
        \mathrm{MHA}\!\left(
          \mathbf{H}^{\mathrm{local}},\,
          \mathbf{H}^{\mathrm{local}},\,
          \mathbf{H}^{\mathrm{local}}
        \right)
      \right)
    \right)
  \label{eq:sa_lcf}
\end{equation}

\noindent trong đó $\mathrm{MHA}(\cdot)$ là phép multi-head attention
với 8 head.

\paragraph{Bước 2: Hợp nhất hai luồng.}
Luồng cục bộ và luồng toàn cục được ghép nối theo chiều ẩn rồi chiếu
xuống $H$ chiều:

\begin{equation}
  \mathbf{H}^{\mathrm{fused}} = \mathrm{Dropout}\!\left(
    \mathbf{W}_2 \left[
      \widetilde{\mathbf{H}}^{\mathrm{local}};
      \mathbf{H}
    \right]
  \right), \quad \mathbf{W}_2 \in \mathbb{R}^{H \times 2H}
  \label{eq:fusion}
\end{equation}

\paragraph{Bước 3: Self-attention tổng hợp và pooling.}
Biểu diễn tổng hợp tiếp tục qua một lớp self-attention $\mathrm{SA}$
nữa trước khi lấy biểu diễn token \texttt{[CLS]}:

\begin{equation}
  \mathbf{H}^{\mathrm{out}} = \mathrm{SA}\!\left(
    \mathbf{H}^{\mathrm{fused}}
  \right)
  \label{eq:sa_final}
\end{equation}

\begin{equation}
  \mathbf{p} = \mathrm{Tanh}\!\left(
    \mathbf{W}_p\, \mathbf{H}^{\mathrm{out}}_{[0]}
  \right), \quad \mathbf{W}_p \in \mathbb{R}^{H \times H}
  \label{eq:pooler}
\end{equation}

\noindent trong đó $\mathbf{H}^{\mathrm{out}}_{[0]}$ là hàng ứng với
vị trí \texttt{[CLS]} (chỉ số 0).

\paragraph{Bước 4: Phân loại đa nhiệm.}
Từ biểu diễn câu $\mathbf{p} \in \mathbb{R}^H$, hai đầu ra được tính
song song:

\begin{align}
  \hat{\mathbf{y}}^{\mathrm{sent}} &= \mathbf{W}_s\, \mathbf{p},
    \quad \mathbf{W}_s \in \mathbb{R}^{3 \times H}
    \label{eq:head_sentiment}\\[4pt]
  \hat{\mathbf{y}}^{\mathrm{cat}} &= \mathbf{W}_c\, \mathbf{p},
    \quad \mathbf{W}_c \in \mathbb{R}^{6 \times H}
    \label{eq:head_category}
\end{align}

\noindent với $\hat{\mathbf{y}}^{\mathrm{sent}}$ là logit cho 3 lớp
cảm xúc (Positive / Negative / Neutral) và
$\hat{\mathbf{y}}^{\mathrm{cat}}$ là logit cho 6 nhóm khía cạnh
(SERVICE, FACILITY, AMENITY, EXPERIENCE, BRANDING, LOYALTY).

Hàm mất mát tổng hợp là tổng của hai hàm cross-entropy có trọng số
lớp:

\begin{equation}
  \mathcal{L} = \mathcal{L}_{\mathrm{CE}}^{\mathrm{sent}}
              + \mathcal{L}_{\mathrm{CE}}^{\mathrm{cat}}
  \label{eq:loss}
\end{equation}

Hình~\ref{fig:two_stream} minh họa toàn bộ kiến trúc hai luồng.

% =====================================================================
\section{Token Merging cho chuỗi 1D}
\label{sec:pp_tome}
% =====================================================================

\subsection{Tổng quan và vị trí tích hợp}
\label{subsec:tome_overview}

Token Merging (ToMe)~\cite{bolya2023tome} được tích hợp vào mô hình
\texttt{FastLcfBertMultiTask} dưới dạng module
\texttt{ToMeSequenceMerger}, áp dụng \textbf{sau} khi backbone BERT
hoàn thành quá trình mã hóa (post-BERT ToMe) và trước khi luồng cục bộ
được tính.

Module nhận đầu vào ba tensor:
$\mathbf{H} \in \mathbb{R}^{B \times L \times H}$ (trạng thái ẩn),
$\mathbf{V} \in \mathbb{R}^{B \times L}$ (véc-tơ chỉ thị khía cạnh)
và $\mathbf{M} \in \{0,1\}^{B \times L}$ (attention mask),
và trả về bộ $(\mathbf{H}', \mathbf{V}', \mathbf{M}')$ có độ dài
chuỗi $L' \leq L$.

Số vòng gộp (merge steps) được cấu hình qua
\texttt{num\_merge\_steps}$= 2$ (mặc định); mỗi vòng áp dụng một
trong ba chiến lược sau.

% ---------------------------------------------------------------------
\subsection{Phép gộp cơ bản và bảo toàn LCF}
\label{subsec:merge_op}
% ---------------------------------------------------------------------

Bất kể chiến lược nào, khi hai token $s$ (source) và $d$ (destination)
được chọn để gộp, phép gộp thực hiện hai thao tác:

\paragraph{Cập nhật biểu diễn bằng trung bình:}
\begin{equation}
  \mathbf{h}'_s = \frac{\mathbf{h}_s + \mathbf{h}_d}{2}
  \label{eq:merge_avg}
\end{equation}

\paragraph{Bảo toàn thông tin khía cạnh bằng max:}
\begin{equation}
  v'_s = \max(v_s,\; v_d)
  \label{eq:merge_lcf}
\end{equation}

Token $d$ bị loại khỏi chuỗi; token $s$ thay thế với giá trị đã được
cập nhật.  Quy tắc \eqref{eq:merge_lcf} đảm bảo rằng nếu một trong
hai token thuộc cụm từ khía cạnh (có $v = 1$), token hợp nhất kết quả
cũng mang nhãn khía cạnh, tránh mất tín hiệu vị trí cho cơ chế LCF
ở bước tiếp theo.

Tất cả ba chiến lược dùng \textbf{độ tương đồng cosine} trên biểu
diễn đã chuẩn hóa $L_2$ để đo mức độ giống nhau giữa hai token:

\begin{equation}
  \hat{\mathbf{h}}_i = \frac{\mathbf{h}_i}{\|\mathbf{h}_i\|_2 + \varepsilon},
  \quad \varepsilon = 10^{-6}
  \label{eq:l2norm}
\end{equation}

\begin{equation}
  \mathrm{sim}(i, j) = \hat{\mathbf{h}}_i \cdot \hat{\mathbf{h}}_j
  \label{eq:cosine}
\end{equation}

% ---------------------------------------------------------------------
\subsection{Chiến lược 1: Bipartite Matching}
\label{subsec:bipartite}
% ---------------------------------------------------------------------

Chiến lược này là chiến lược gốc của ToMe~\cite{bolya2023tome}, được
điều chỉnh cho chuỗi 1D và bổ sung bảo vệ token khía cạnh.

\paragraph{Xây dựng tập ứng viên.}
Từ $n$ token hợp lệ của câu (không tính padding), tập ứng viên gộp
$\mathcal{P}$ được xác định bằng cách loại trừ:
\begin{itemize}
  \item Token \texttt{[CLS]} ($i = 0$), bảo vệ bởi \texttt{protect\_cls};
  \item Token \texttt{[SEP]} cuối ($i = n-1$), bảo vệ bởi
    \texttt{protect\_sep};
  \item Các token khía cạnh ($v_i > 0.5$), bảo vệ bởi
    \texttt{protect\_aspect}.
\end{itemize}

\paragraph{Phân chia hai nhóm.}
Các phần tử của $\mathcal{P}$ được liệt kê theo thứ tự chỉ số tăng
dần thành dãy $p_0, p_1, \dots, p_{|\mathcal{P}|-1}$.  Dãy này được
chia xen kẽ thành hai nhóm:

\begin{align}
  \mathcal{A} &= \{p_0,\, p_2,\, p_4, \dots\} \quad \text{(chỉ số
    chẵn trong dãy)}\label{eq:group_A}\\
  \mathcal{B} &= \{p_1,\, p_3,\, p_5, \dots\} \quad \text{(chỉ số
    lẻ trong dãy)}\label{eq:group_B}
\end{align}

\paragraph{Tính ma trận độ tương đồng.}
Mỗi hàng ứng với một token trong $\mathcal{A}$, mỗi cột ứng với một
token trong $\mathcal{B}$:

\begin{equation}
  \mathbf{S} = \widehat{\mathbf{A}}\, \widehat{\mathbf{B}}^{\top},
  \quad \mathbf{S} \in \mathbb{R}^{|\mathcal{A}| \times |\mathcal{B}|}
  \label{eq:bipartite_score}
\end{equation}

\noindent trong đó $\widehat{\mathbf{A}}, \widehat{\mathbf{B}}$ là ma
trận biểu diễn chuẩn hóa $L_2$ của các token trong nhóm tương ứng.

\paragraph{Ghép cặp tham lam.}
Mỗi token $a_i \in \mathcal{A}$ chọn token $b_{j^*} \in \mathcal{B}$
có độ tương đồng lớn nhất:

\begin{equation}
  j^*(i) = \arg\max_{j} S_{ij}
  \label{eq:greedy_assign}
\end{equation}

Ràng buộc một-một phía $\mathcal{B}$: mỗi $b_j$ chỉ được gộp với một
$a_i$ duy nhất (cặp đầu tiên yêu cầu $b_j$ được giữ lại, các cặp sau
yêu cầu cùng $b_j$ bị bỏ qua).  Đây là phép ghép cặp hai phía
\textit{greedy} theo nghĩa mỗi bên $\mathcal{A}$ giải quyết độc lập
mà không có giai đoạn đấu giá toàn cục.

Sau bước này, mỗi cặp $(a_i, b_{j^*(i)})$ được gộp theo
\eqref{eq:merge_avg}--\eqref{eq:merge_lcf} với $s = a_i$, $d =
b_{j^*(i)}$.

\paragraph{Độ phức tạp.}
Với $|\mathcal{P}|$ token ứng viên và chiều ẩn $H$, chi phí tính ma
trận $\mathbf{S}$ là $\mathcal{O}(|\mathcal{P}|^2 / 4 \cdot H)$, thấp
hơn nhiều so với self-attention $\mathcal{O}(L^2 H)$ trên toàn chuỗi.

% ---------------------------------------------------------------------
\subsection{Chiến lược 2: Sequential Local Merging}
\label{subsec:sequential}
% ---------------------------------------------------------------------

Chiến lược này thực hiện gộp \textbf{tuần tự từ trái sang phải}, mỗi
token quyết định gộp dựa trên so sánh với hai lân cận gần nhất.

\paragraph{Xây dựng tập bảo vệ.}
Ngoài CLS, SEP và token khía cạnh, chiến lược này bổ sung bảo vệ
\textbf{ranh giới mid-SEP}: trong định dạng SPC
(xem~\eqref{eq:spc}), token phân cách giữa phần câu văn và phần
khía cạnh (tức \texttt{[SEP]} đầu tiên, không phải token cuối) không
được xóa.  Điều kiện nhận biết:

\begin{equation}
  \mathrm{midSep}(j) = \mathbf{1}\!\left[
    v_j < 0.5 \;\text{ và }\; v_{j+1} > 0.5
  \right]
  \label{eq:midsep}
\end{equation}

\paragraph{Vòng quét.}
Biến con trỏ $i$ duyệt từ $\texttt{protect\_left}$ đến $n -
\texttt{protect\_right} - 1$.  Tại mỗi bước $i$:

\begin{enumerate}
  \item Bỏ qua nếu $i$ đã bị xóa, là mid-SEP hoặc là token khía
    cạnh.
  \item Tìm \textbf{lân cận trái tích cực gần nhất}:
    $l = \max\{l' < i : l' \text{ chưa xóa và không là mid-SEP}\}$.
  \item Tìm \textbf{lân cận phải tích cực gần nhất không phải
    khía cạnh}:
    $r = \min\{r' > i : r' \text{ chưa xóa, không mid-SEP, không
    aspect}\}$.
  \item Tính độ tương đồng:
    \begin{align}
      \mathrm{sim}_{\mathrm{L}} &=
        \mathrm{sim}(i, l) \quad \text{(nếu tồn tại } l \text{)}
        \label{eq:sim_left}\\
      \mathrm{sim}_{\mathrm{R}} &=
        \mathrm{sim}(i, r) \quad \text{(nếu tồn tại } r \text{)}
        \label{eq:sim_right}
    \end{align}
  \item Ra quyết định gộp:
    \begin{equation}
      \begin{cases}
        \text{Token } i \text{ gộp vào trái}:
          & \mathbf{h}_l \leftarrow \tfrac{\mathbf{h}_l + \mathbf{h}_i}{2},\;
            v_l \leftarrow \max(v_l, v_i),\;
            \text{xóa } i \\
            & \quad \text{khi } \exists l \text{ và }
              \mathrm{sim}_{\mathrm{L}} > \mathrm{sim}_{\mathrm{R}}\\[6pt]
        \text{Token phải gộp vào } i:
          & \mathbf{h}_i \leftarrow \tfrac{\mathbf{h}_i + \mathbf{h}_r}{2},\;
            v_i \leftarrow \max(v_i, v_r),\;
            \text{xóa } r \\
            & \quad \text{khi } \exists r \text{ và }
              \mathrm{sim}_{\mathrm{R}} \geq \mathrm{sim}_{\mathrm{L}}
      \end{cases}
      \label{eq:seq_decision}
    \end{equation}
  \item Tăng $i \leftarrow i + 1$.
\end{enumerate}

Một lần duyệt đầy đủ ($i$ từ đầu đến cuối) được gọi là một
\textit{merge step}; mô hình thực hiện \texttt{num\_merge\_steps}$= 2$
lần duyệt liên tiếp.

\paragraph{Tính chất.}
Token có thể \textit{nhận} nhiều gộp từ nhiều hướng trong cùng một
lần duyệt (không có ràng buộc một-một phía đích).  Điều này khác với
Bipartite Matching và cho phép co chuỗi nhanh hơn khi có nhiều cụm
token lặp lại liên tiếp.

% ---------------------------------------------------------------------
\subsection{Chiến lược 3: Attention-Weighted Merging}
\label{subsec:attn_merge}
% ---------------------------------------------------------------------

Chiến lược này gộp ưu tiên \textbf{các token có trọng số chú ý thấp
nhất}, đồng thời bảo vệ các token mang thông tin quan trọng.

\paragraph{Xây dựng tập bảo vệ cứng.}
Ba nhóm token không bao giờ bị xóa:

\begin{itemize}
  \item \textbf{Biên cấu trúc}: CLS ($i = 0$) và SEP ($i = n-1$);
  \item \textbf{Token khía cạnh}: $v_i > 0.5$;
  \item \textbf{Token chú ý cao}: những token có trọng số chú ý
    $a_i$ nằm trong \textbf{phân vị thứ 75} trở lên:
    \begin{equation}
      \theta = \mathrm{quantile}(\mathbf{a},\; 0.75)
      \label{eq:attn_threshold}
    \end{equation}
    \begin{equation}
      \mathrm{highAttn}(i) = \mathbf{1}[a_i \geq \theta]
      \label{eq:high_attn}
    \end{equation}
\end{itemize}

Token $i$ bị bảo vệ khi thỏa ít nhất một trong các điều kiện trên.

\paragraph{Giới hạn số cặp gộp.}
\begin{equation}
  T_{\max} = \left\lfloor
    \frac{n - \texttt{protect\_left} - \texttt{protect\_right}}{2}
  \right\rfloor
  \label{eq:max_merges}
\end{equation}

\paragraph{Vòng lặp gộp lặp.}
Trong mỗi bước $t = 1, 2, \dots, T_{\max}$:

\begin{enumerate}
  \item \textbf{Tìm ứng viên bị xóa}: token chỉ số $i$ có trọng số
    chú ý nhỏ nhất trong số các token chưa bị xóa và không thuộc tập
    bảo vệ:
    \begin{equation}
      i = \arg\min_{\substack{t:\, \text{not removed}(t)\\
            \text{not protected}(t),\, v_t \leq 0.5}}
        a_t
      \label{eq:find_min_attn}
    \end{equation}
    Nếu không có ứng viên: dừng vòng lặp sớm.

  \item \textbf{Tìm token đích}: token $j \neq i$ chưa bị xóa, không
    phải khía cạnh, có độ tương đồng cosine với $i$ cao nhất:
    \begin{equation}
      j = \arg\max_{\substack{k \neq i:\, \text{not removed}(k)\\
            v_k \leq 0.5}}
        \hat{\mathbf{h}}_i \cdot \hat{\mathbf{h}}_k
      \label{eq:find_neighbor}
    \end{equation}
    Nếu giá trị tương đồng $< -1.5$ (không có lân cận hợp lệ): dừng.

  \item \textbf{Thực hiện gộp}: gộp $i$ vào $j$ theo
    \eqref{eq:merge_avg}--\eqref{eq:merge_lcf}, đánh dấu $i$ đã xóa.

  \item \textbf{Cập nhật metric}: chuẩn hóa lại toàn bộ ma trận biểu
    diễn trước vòng lặp kế tiếp.
\end{enumerate}

\paragraph{Sự khác biệt chính so với hai chiến lược còn lại.}
Bipartite Matching và Sequential Local chọn cặp dựa hoàn toàn vào độ
tương đồng cosine giữa biểu diễn.  Attention-Weighted Merging sử dụng
trọng số chú ý $\mathbf{a}$ như một \textit{tín hiệu tầm quan trọng}
độc lập để ưu tiên xóa các token mang ít thông tin cảm xúc nhất, đồng
thời bảo vệ không chỉ token khía cạnh mà cả những token có cơ chế chú
ý của BERT đã học được rằng chúng quan trọng.

% ---------------------------------------------------------------------
\subsection{Chế độ resize và compact}
\label{subsec:resize}
% ---------------------------------------------------------------------

Sau tất cả các vòng gộp, chuỗi có độ dài $L' \leq L$.  Module
\texttt{ToMeSequenceMerger} hỗ trợ hai chế độ đầu ra, điều khiển bởi
tham số \texttt{resize}:

\paragraph{Chế độ \texttt{resize=True} (nội suy về độ dài gốc).}
Chuỗi $\mathbf{H}' \in \mathbb{R}^{L' \times H}$ được nội suy tuyến
tính về đúng độ dài $L$ ban đầu:

\begin{equation}
  \mathbf{H}^{\mathrm{out}} =
    \mathrm{Interpolate}\!\left(\mathbf{H}',\; L,\;
    \text{mode=linear}\right)
    \in \mathbb{R}^{L \times H}
  \label{eq:resize}
\end{equation}

Chế độ này \textbf{không thay đổi hình dạng tensor} nên các lớp hạ
nguồn (SA, Linear) không cần điều chỉnh.  Nó cô lập tác động của
gộp token lên \textit{chất lượng biểu diễn} mà không đo tốc độ suy
luận.  Đây là chế độ được dùng trong toàn bộ các thực nghiệm chính của
luận văn.

\paragraph{Chế độ \texttt{resize=False} (compact).}
Chuỗi rút gọn $\mathbf{H}'$ được giữ nguyên; attention mask mới
$\mathbf{M}' \in \{0, 1\}^{B \times L'_{\max}}$ được tính và truyền
xuống các lớp self-attention.  Chế độ này mang lại \textbf{tăng tốc
thực sự} ($\mathcal{O}({L'}^2)$ thay cho $\mathcal{O}(L^2)$) nhưng
yêu cầu các lớp hạ nguồn hỗ trợ chuỗi có độ dài thay đổi.  Đây là
hướng tối ưu hóa được đề xuất cho các nghiên cứu tiếp theo.

\paragraph{Tóm tắt hai chế độ.}

\begin{table}[h]
  \centering
  \caption{So sánh hai chế độ đầu ra của \texttt{ToMeSequenceMerger}.}
  \label{tab:resize_mode}
  \begin{tabular}{lcc}
    \toprule
    \textbf{Thuộc tính} & \texttt{resize=True} & \texttt{resize=False} \\
    \midrule
    Độ dài đầu ra & $L$ (giữ nguyên) & $L' \leq L$ \\
    Hình dạng tensor thay đổi & Không & Có \\
    Tăng tốc suy luận thực sự & Không & Có \\
    Mục tiêu đánh giá & Chất lượng biểu diễn & Tốc độ + chất lượng \\
    Dùng trong thực nghiệm chính & \checkmark & \\
    \bottomrule
  \end{tabular}
\end{table}
